\chapter{线上拓展专题}
\label{online:topics}

本页内容不印入纸质教材，是正文拓展层的延伸阅读，供课程设计、毕业设计和自学使用。各节开头注明对应的正文位置；先完成正文相应小节，再阅读这里的展开内容。

\section{数字孪生平台的工程化细节}\label{online:twin}

\paragraph{对应正文}第6章6.4节（数字孪生水利平台架构概览）与第8章8.5节（数字孪生水利平台案例）。正文保留概念、分层、数据与模型、闭环与契约、典型场景和运行值守的主线；本节收录跨层契约的完整论述、验收检查项、课程项目迭代和发展方向。

\subsection{跨层契约与架构演进}

每层的接口文档都应写明输入、输出、失败响应与责任主体。物理对象层提供资产目录、设备状态和工程边界；感知连接层提供带事件时间的原始消息与接入证据；数据底板层提供经过质量码标记的可查询数据产品；模型与知识层提供带适用范围和不确定性的计算结果；服务层提供查询、预警、预演和调度建议；交互应用层提供证据浏览、人工确认和操作回执。任一层缺少版本或责任字段，上一层的正确性都无法传递到下一层。

契约应同时覆盖正常流和失败流。正常流中，网关接收带有对象编码和事件标识的观测，质量服务完成单位、时间和范围校验，数据底板写入不可变原始记录，模型服务读取符合条件的数据快照，业务服务生成带证据的预警或方案，应用层显示并等待授权操作。失败流中，缺少对象编码的消息进入隔离队列，单位错误被标记为可疑，关键测点缺测使模型降级或停算，模型超时返回可解释的不可用状态，权限不足只返回允许范围内的数据。每一种失败都要有稳定错误码、责任角色和恢复动作，不能把异常堆栈直接当作业务提示。

架构评审可采用“数据向上、指令向下、证据贯穿”的三条线。数据向上检查采集、质检、存储和模型输入是否保持同一对象编码、时间基准和质量语义；指令向下检查建议、审批、执行和回执是否经过权限、有效期和设备状态校验；证据线横向关联原始观测、模型运行、规则版本、界面操作和工单结果。三条线交汇处就是平台的审计边界：例如某次闸门调度建议既要能追到降雨情景和预报运行，也要能追到审批人、设备回执和执行后的水位变化。

平台演进应优先扩展契约，再替换实现。把单工程场景扩展到流域时，新增的是河网拓扑、断面、预报模型和影响对象目录，原有测点、质量码、事件时间和权限语义仍然有效；把本地三维场景扩展到地球级浏览时，增加的是大范围坐标、瓦片服务和LOD策略，业务对象编码和时间窗不能因切换引擎而改变。若某个新模块要求修改既有字段含义，应发布新的契约版本，同时保留一段兼容期并记录迁移规则；直接复用旧字段承载新语义，会让历史数据和新结果无法比较。

模型、数据和前端的发布节奏也应解耦。模型服务可以先影子运行，使用真实输入生成不影响生产的结果；数据底板可以先扩展字段，旧客户端继续读取兼容视图；前端可以先增加“模拟”标签和版本显示，等待后端接口正式开放。发布单要列出依赖关系、回滚条件和观测指标，运维员依据指标逐步放量。发生异常时，先停止新版本进入决策链，再保留已经产生的预警和工单，最后按输入快照和版本号重放，判断异常来自数据、模型还是界面。

跨层契约还决定课程项目的分工方式。负责三维场景的学生提交对象编码、坐标转换和加载错误证据；负责数据服务的学生提交质量检查、时序查询和幂等写入证据；负责模型的学生提交输入快照、模型卡和不确定性说明；负责业务闭环的学生提交预警、工单、审批和回执证据。各组以同一组样例数据和同一追踪号联调，任何一组都不能用手工复制的“演示数字”替代接口调用。接口检查通过后，各组再组合运行；出现失败时，按追踪号核对相关层的输入与输出，确定需要修改的模块。

架构设计还要明确哪些状态可以自动推进，哪些状态必须等待人工确认。数据接入和质量标记属于可自动化的机械步骤，但“是否进入风险评分”应由质量规则和模型契约共同决定；模型运行可以自动排队和重试，但模型结果进入预警或调度建议前要检查适用范围和输入快照；预警事件可以自动生成，但高风险处置、闸门命令和预案启用必须由授权角色审批。把自动步骤和人工步骤写成状态机后，系统才能在界面上给出“待复核、已批准、执行中、待回执”等明确状态，而不是用一个“完成”字段掩盖责任交接。

跨层接口的时间语义尤其容易被忽略。采集设备的事件时间用于重建现场事实，网关接收时间用于衡量链路延迟，模型开始和结束时间用于衡量计算资源，人工确认时间用于衡量处置时效。四类时间都要保留，不能把接收时间覆盖事件时间，也不能把页面刷新时间当成模型结果时刻。发生迟到数据时，系统按事件时间判断是否需要重算；发生模型超时时，系统按处理时间触发降级；发生审批超时时，系统按业务有效期升级值守任务。时间字段的分工一旦写入契约，跨服务日志才具有可比较性。

空间语义也应沿层传递。物理对象层保存工程坐标和对象关系，数据底板同时保存水平CRS、高程基准和转换版本，模型输入输出保留计算网格或断面坐标，服务层返回对象编码与空间范围，应用层才把它们转换成局部原点、地图瓦片或三维屏幕位置。任何一层改变坐标，都必须在结果中写出转换记录；如果只在前端拖动模型使其“看起来对齐”，后端查询、剖面分析和监测点回放仍会使用错误位置。架构验收应选择至少一个测点和一个工程构件，依次核对原始坐标、标准坐标、局部坐标、三维拾取及反向查询的结果。

当平台需要接入新模型或新数据源时，先做契约评审，再做适配器开发。评审清单包括对象编码映射、时间基准、单位和质量码、输入输出版本、错误与超时、权限和审计、性能预算以及回滚方案。适配器负责把外部格式转换为平台规范，不把外部字段直接泄露到业务页面；对于无法映射的字段，进入扩展属性并在数据字典登记。这样，模型可以替换、数据源可以增加、三维引擎可以升级，而既有预警、工单和审计记录仍能按原语义读取。

最终的架构验收以业务证据为中心组织。抽取一条观测事件，验证它能从接入日志追到质量结论、状态版本、模型输入、风险结果、人工确认、指令回执和复盘记录；再制造一条重复消息、一个单位错误、一次模型超时和一次无权限命令，验证系统分别给出幂等忽略、可疑标记、降级结果和拒绝审计。正常用例检查输出是否符合预期，失败用例检查处理停止的位置、恢复步骤及对应责任人。第8章的案例实现沿用这套证据链，读者可在\ref{sec:ch08-twin-acceptance}节看到端到端脚本和阶段门禁的具体写法。

可以用一次“连续降雨后渗压变化”的演算检查双向闭环：雨量事件上行后，质量服务标记时间窗和数据完整性，模型服务根据水位与渗压快照计算趋势，预警服务给出带可信状态的结果；值班员确认后生成工单，若需要调度则由审批人选择方案并下发命令，设备回执和新的水位、流量又沿数据链上行。每一步都保留追踪号和版本，才能在复盘时回答“当时看到了什么、依据哪一版模型、谁作了决定、执行后发生了什么”。这组记录对应连接维的数据与反馈过程，可用于检查模型建议、人工决策和设备执行是否一致。
联调时还应验证这一路径的失败分支：人为制造迟到观测、模型超时和审批拒绝，检查状态是否停在正确节点，并核对恢复后是否补齐证据。通过同一追踪号比较正常与失败两条路径，学生能够把抽象的五维关系落实为可观察的系统行为。
验收报告应附状态转移、输入输出快照和责任签名，便于跨角色复核；复核记录写出模型版本、数据窗口、异常处置人、复核时间与责任角色。责任角色未签名时，验收状态保持“待复核”，不进入正式发布；状态变更、签名和发布动作使用同一追踪号关联，并由审计角色独立复核。证据链完整，系统结果才可进入正式业务处置。

\subsection{各小节的补充要求}

\paragraph{模型变更流程}
模型变更遵循“开发—离线验证—影子运行—评审—发布—监控—回滚”流程。影子运行让新模型接收真实输入但不影响生产决策，可比较新旧结果并发现边界工况。模型审批记录与软件发布记录共同构成审计证据。

\paragraph{单位与时区}
单位转换在数据标准层集中完成。例如水位统一为米、流量统一为立方米每秒；接口仍携带单位字段。阈值和模型参数必须声明所用单位，禁止依赖开发者记忆。时区统一存储为带偏移的时间或UTC，界面按用户时区显示。

\paragraph{三维交互的补充要求}
三维界面应帮助用户定位对象、理解状态和获取证据。对象选择后显示稳定业务名称、编码、数据时刻和质量状态；颜色表达必须配合图例、文字或图标，不能只用红绿色区分风险。对键盘操作、字号和对比度也应进行基本无障碍检查。

时间轴允许回看任一历史状态，但必须区分“当时已知的数据”和“后来补录后重算的状态”。若回放使用了后来数据，界面应标注为修订视图，防止把事后计算结果误认为当时决策依据。

空间查询同样需要证据边界。淹没范围、影响人口或工程数量由特定模型和数据版本计算，界面点击统计数字时应能打开其空间范围、过滤条件和生成时间。截图只能作为沟通材料，正式结果应保留可查询数据和生成记录。

三维性能下降时优先保证业务信息可达。模型可降级LOD或切换二维地图，告警列表、数据表和处置入口仍应可用。通过这种渐进增强，三维场景不会成为整个业务系统的单点故障。

\paragraph{契约测试}
接口契约需通过自动化测试验证，包括必填字段、单位、时区、枚举、重复事件和乱序事件。契约文档与代码同步版本化，避免不同团队凭口头约定解释同一字段。

\paragraph{方案比选的展示要求}
方案比选至少展示目标、约束、模型版本、边界条件、关键断面过程和风险指标。任何调度建议都要保留输入快照；预报更新后重新计算时，新旧方案不能混用同一编号。

表中“反馈”环节将调度执行后的实测水位和流量送入状态更新，并保留原预报结果供误差比较；若执行回执缺失，系统只能把方案标记为“已批准、待回执”，不能假设闸门已经按计划动作。预报模型的误差评价关注洪峰、峰现时间和过程拟合，调度模型的评价则关注约束满足、供水缺口、生态下泄和下游风险，两类指标分开统计后才能判断是预报误差还是方案执行偏差。

\paragraph{部署设计要点}
部署设计需明确网络分区、离线时长、数据补传顺序、模型降级和灾备目标。例如模型服务不可用时，平台仍应展示经质检的实时数据和已批准阈值规则；不能把整个监控能力绑定在单一模型上。

\paragraph{安全测试}
高风险控制采用职责分离：分析人员提出建议，授权人员审批，执行系统校验有效期和设备状态。系统不得因“自动化”跳过工程运行规程。紧急手动操作也要在恢复后补录原因和结果。

安全测试包括越权访问、令牌失效、重放、消息篡改、接口限流和审计完整性。三维前端同样属于攻击面，模型文件、纹理和属性数据要进行来源校验，避免把不可信脚本或敏感属性直接交付浏览器。

\paragraph{可观测性与值守的补充要求}
平台上线后，运行团队必须回答三个问题：当前服务是否可用，当前数据是否可信，当前模型结果能否用于业务。仅监控CPU和内存无法回答后两个问题，因此需要把服务、数据、模型和业务闭环统一纳入可观测体系。表\ref{tab:dt-observability}从服务、数据、模型、业务和安全五方面列出指标，分别检查服务可用性、数据新鲜度与完整性、模型适用范围、预警处置进度，以及权限与审计异常。CPU和内存正常时，仍可能发生断报、模型越界或工单积压，需要逐类观察。每条关键链路应配置服务级目标（Service Level Objective，SLO），并明确统计窗口、允许失败比例、告警责任人和处置时限。

数据血缘记录数据集从来源到结果的变换关系。每个血缘节点应包含来源标识、时间范围、处理规则版本、输入与输出校验摘要、执行状态和责任主体；模型结果还要关联模型卡、参数集和运行环境。采集、计算和发布流程应自动将血缘记录写入元数据存储，并能从三维对象或告警详情反向查询。

值守界面应按影响范围和紧迫程度合并告警，避免同一根因产生大量重复提示。告警必须包含对象、发生时间、当前值、质量码、触发规则、推荐动作和确认入口；抑制、合并或关闭告警也要保留操作者与理由。交接班记录尚未恢复的数据源、降级中的模型、未闭环工单和临时权限，防止信息只停留在即时通信中。

当数据质量或模型表现持续低于阈值时，系统应自动降低结果可信等级，必要时退出自动推荐，但仍保留原始数据、二维地图和人工规程。恢复正常不能只看服务重新启动，还应完成积压补传、时间顺序校验、状态重算、告警去重和责任人确认。值守人员核对这些恢复记录后，才能确认数据、模型与业务处置已恢复到可用状态。

\subsection{验证、验收与持续改进}\label{online:twin-acceptance}

验收需要将数据、空间、模型、服务、反馈过程和运维指标对应到具体输入、版本、责任人与复核结果。案例水库的失效降级步骤、验收脚本和阶段门禁见\ref{sec:ch08-twin-acceptance}节，可用于检查各项结果是否能够重复获得。

数字孪生验收既包含软件质量，也包含数据和模型可信度。建议分为数据、空间、模型、服务、闭环和运维六类指标。表\ref{tab:dt-acceptance}列出这六类的具体检查项。每项结果都应附可复查的记录。例如模型精度检查需列出检查数据、适用工况、误差指标及其与验收限值的比较。

\begin{table}[htbp]
\centering
\caption{数字孪生平台验收检查项}
\label{tab:dt-acceptance}
\begin{tabular}{p{2.3cm}p{5.0cm}p{5.4cm}}
\toprule
类别 & 检查项 & 证据 \\
\midrule
数据 & 完整率、时效、单位、质量码、血缘 & 抽样记录与异常处置日志 \\
空间 & CRS、高程基准、控制点、模型对齐 & 校核报告与同名点误差 \\
模型 & 适用范围、参数、校准、回放结果 & 模型卡、版本和测试数据 \\
服务 & 可用性、延迟、错误处理、权限 & 压测、安全测试和监控 \\
闭环 & 告警到工单、指令到回执、责任追踪 & 演练记录与审计链 \\
运维 & 备份、降级、灾备、变更和回滚 & 运维手册与恢复演练 \\
\bottomrule
\end{tabular}
\end{table}

上线前使用历史典型洪水、设备故障和缺测场景回放；上线后持续统计模型残差、告警准确性、处置时长和数据质量。改进必须经过版本化、测试、审批和回滚准备。

\subsection{课程项目实施路线}\label{online:twin-course}

课程项目按对象与坐标、场景绑定、模型服务、业务处置的顺序推进，每轮提交代码、数据样例与运行记录。案例水库的端到端实现、验收脚本和分阶段交付见\ref{sec:ch08-twin-course}节，可用于核对各轮成果的接口和验证步骤。

课程项目不要求一次实现完整流域数字孪生，可分四次迭代，每次形成可验收成果。表\ref{tab:dt-course-iterations}把四轮的任务、交付物和验收重点逐行列开。各轮成果依次提供后续实现的输入：对象编码和数据字典支撑三维定位，空间与观测状态为模型提供输入，模型事件再进入处置流程。选用经验公式时也需记录其适用条件，并验证输入与结果。第一轮建立对象编码、空间基准和数据样例；第二轮完成三维场景与状态绑定；第三轮接入一个模型或规则并形成事件；第四轮完成处置闭环、测试和演示。

\begin{table}[htbp]
\centering
\caption{数字孪生课程项目的四次迭代}
\label{tab:dt-course-iterations}
\begin{tabular}{p{1.8cm}p{4.0cm}p{4.1cm}p{3.3cm}}
\toprule
迭代 & 主要任务 & 交付物 & 验收重点 \\
\midrule
一 & 场景选择、对象编码、CRS与数据字典 & 范围说明、数据样例 & 边界和元数据完整 \\
二 & 加载GIS/三维模型、绑定对象状态 & 可交互页面、映射表 & 定位正确、异常可见 \\
三 & 实现质检、规则或简化模型 & 模型卡、结果事件 & 输入契约、版本可追踪 \\
四 & 告警处置、审计、测试和降级 & 演示系统、验收报告 & 闭环可复现、失败可降级 \\
\bottomrule
\end{tabular}
\end{table}

团队角色可分为数据/空间、前端场景、后端服务和测试集成，但接口与对象编码必须共同评审。每次迭代使用同一组场景用例，避免各成员分别制作无法拼接的演示。

最终演示应包含正常流、数据质量异常、模型不可用和无权限操作四类场景。验收报告记录输入文件、运行版本、操作步骤、预期和实际结果，使其他小组能够复现，而不仅是录制一段成功视频。

\subsection{发展重点}

数字孪生的持续改进可从模型可信度、语义互操作、云边协同、事件处理、人工决策支持和安全六方面展开，并分别设置验证指标。

模型可信度与不确定性表达，衡量的是模型卡是否记录了适用范围与已知失效条件，以及预报结果能否给出区间而不只是一个数；验证方式是用历史场景回放，统计实测值落在给出区间内的比例。跨系统语义互操作，衡量的是同一座水库在不同系统里的编码、名称和特征水位是否一致；验证方式是抽取若干工程做跨库比对，冲突数应当趋近于零。云边协同，衡量的是断网期间边缘侧能否继续采集与本地判断、恢复后数据是否按事件时间正确回补；验证方式是拉断网络若干分钟再接上，检查有无缺测空洞和重复入库。

面向事件的实时数据处理，衡量的是从观测到达到预警生成的端到端时延，以及消息积压在汛期峰值下的恢复速度。人在回路的决策支持，衡量的是每一条进入执行的指令是否都能回溯到审批人、依据的模型版本和输入快照。网络与数据安全，衡量的是权限越界尝试能否被拒绝并留痕、密钥能否在不停机的情况下轮换。

将上述方向列入建设方案时，还应给出指标的统计窗口、测试数据和失败条件，以便比较实施前后的效果。

\subsection{底板与算法解耦}

可复制的平台把空间底板、业务数据、模型算法和应用编排分开。更换洪水模型时，不应重新制作工程模型和用户体系；新增工程时，通过元数据和适配器接入，而不是复制一套源码。

图\ref{fig:twin-decoupling}把四者的解耦关系画了出来。检验解耦是否成立有个简单办法：设想更换洪水模型，数一数需要改动的模块——如果空间底板和用户体系也要跟着改，说明分层只停留在文档里。

\begin{figure}[htbp]
\centering
\begin{tikzpicture}[node distance=5mm,
  twinlayer/.style={draw=Primary,rounded corners,fill=PrimaryLight,minimum width=11cm,
    minimum height=.8cm,align=center}]
\node[twinlayer] (app) {业务应用：防洪预演、水网调度、工程安全、会商};
\node[twinlayer,below=of app] (orchestrate) {场景编排：任务、权限、状态机、证据链};
\node[twinlayer,below=of orchestrate] (model) {模型服务：水文、水动力、调度优化、空间分析};
\node[twinlayer,below=of model] (base) {数据与空间底板：对象编码、时空数据、三维模型、元数据};
\node[twinlayer,below=of base] (resource) {算力与边缘资源：调度、监控、缓存、灾备};
\foreach \a/\b in {app/orchestrate,orchestrate/model,model/base,base/resource}{\draw[<->,draw=Primary,thick] (\a)--(\b);}
\end{tikzpicture}
\caption{数字孪生水利平台的解耦分层}
\label{fig:twin-decoupling}
\end{figure}

\subsection{数字孪生平台的技术路线对比}

\paragraph{对应正文}第8章8.5.4节（数字孪生水库矩阵与“四预”）。

数字孪生平台可以从不同的工程起点形成技术路线。以厂商手册为参考的产品集成路线，通常先把流域、水网和工程入口组织成统一门户，再把模型运行、预演任务和工单闭环接到同一套权限与审计机制中。这种路线的优点是业务人员容易按照场景进入系统，适合在已有平台上快速建立可演示、可操作的流程；它的边界是产品中的对象编码、数据质量规则和模型适用范围仍需项目团队重新核验，手册中的示例数据也不能直接替代真实工程数据。

以公开标准和数据底板为中心的路线，则先确定对象目录、时空基准、数据交换协议、元数据和版本规则，再把不同厂商或不同专业模型作为可替换服务接入。它有利于长期积累和跨系统交换，但前期需要投入更多的目录治理、接口设计和质量校验工作。

以数字孪生流域或区域治理为牵引的路线，强调跨工程的联合预报、影响分析和调度协同，重点解决数据范围大、模型链条长和多部门协作的问题；其验收不能只看单个三维页面，而要检查情景输入、模型版本、结果证据和处置回执是否能够闭环。

项目定制路线则围绕某一类水库、水闸或河段的规程和业务动作构建，能够精确贴合现场，但需要用清晰的扩展点避免把一次性约定固化为全行业标准。

四类路线可以组合使用：厂商方案是其中一种实现的参考，公开政策和行业样本提供业务边界，标准化底板提供交换约束，项目规则负责把模型结果转成可审批、可追溯的行动。教学案例采用这种分层比较方式，目的在于训练方案选择和证据核验能力，而不是评价某个产品的优劣。

\section{生产环境的认证与安全细节}\label{online:security}

\paragraph{对应正文}第5章5.6节。正文实现DUTY、ANALYST、OPS三角色和访问令牌；本专题进一步讨论细粒度权限、刷新令牌、撤销与密钥轮换。以下为工程扩展示例，未接入配套最小应用；\texttt{lesson56/RefreshTokenVerifier}单独提供可运行校验器及测试。完整平台仍以正文与配套文件为准。

\subsection{细粒度权限、刷新与撤销的扩展实现}

图\ref{fig:online-ch05-auth-sequence}展示增加刷新端点后的认证流程，区别于正文的重新登录方案。

\begin{figure}[htbp]
\centering
\begin{tikzpicture}[>=Stealth,thick,
  auth/.style={draw=Primary,rounded corners=2pt,fill=PrimaryLight,text=PrimaryDark,text width=2.05cm,minimum height=.9cm,align=center,font=\scriptsize},
  reject/.style={auth,draw=Accent,fill=white}]
  \node[auth] (login) at (0,0) {登录\\提交凭据};
  \node[auth] (issue) at (2.75,0) {认证服务\\签发令牌对};
  \node[auth] (api) at (5.5,0) {API请求\\Bearer access};
  \node[auth] (filter) at (8.25,0) {JWT过滤器\\验签\\类型、有效期};
  \node[auth] (decision) at (11,0) {URL/方法授权\\角色与权限};
  \node[auth] (service) at (13.75,0) {业务服务\\执行用例};
  \node[auth] (refresh) at (2.75,-2) {刷新端点\\验证refresh并轮换};
  \node[reject] (unauth) at (8.25,-2) {401\\认证失败};
  \node[reject] (denied) at (11,-2) {403\\权限不足};
  \foreach \a/\b in {login/issue,issue/api,api/filter,filter/decision,decision/service}
    \draw[->,Primary] (\a)--(\b);
  \draw[->,Accent] (filter)--(unauth);
  \draw[->,Accent] (decision)--(denied);
  \draw[->,Accent,dashed] (unauth.west)--node[above,align=center,font=\scriptsize]{access过期且\\持有有效refresh}(refresh.east);
  \draw[->,Primary] (refresh)--(issue);
  \node[below=.3cm of denied,font=\scriptsize,text=TextGray] {403不触发令牌刷新};
\end{tikzpicture}
\caption{Spring Security 6 与 JWT 的认证授权时序}
\label{fig:online-ch05-auth-sequence}
\end{figure}




清单\ref{lst:online-ch05-security-chain}是图\ref{fig:ch05-security-filter}中那条过滤链的配置：会话置为无状态，认证端点放行，其余请求交给 JWT 过滤器。

\begin{lstlisting}[label={lst:online-ch05-security-chain},language=Java,caption={Spring Security过滤链}]
import org.springframework.context.annotation.*;
import org.springframework.http.HttpStatus;
import org.springframework.security.config.annotation.method.configuration.*;
import org.springframework.security.config.annotation.web.builders.*;
import org.springframework.security.config.http.SessionCreationPolicy;
import org.springframework.security.web.*;
import org.springframework.security.web.authentication.*;

@Configuration
@EnableMethodSecurity
@EnableConfigurationProperties(JwtProperties.class)
class SecurityConfig {
    @Bean
    SecurityFilterChain securityFilterChain(HttpSecurity http,
            JwtAuthenticationFilter jwtFilter) throws Exception {
        return http
            .csrf(AbstractHttpConfigurer::disable)
            // 接入下文的 corsConfigurationSource Bean；
            // 只定义 Bean 而不调用 .cors() 时该配置不会生效
            .cors(Customizer.withDefaults())
            .sessionManagement(s -> s.sessionCreationPolicy(
                    SessionCreationPolicy.STATELESS))
            .exceptionHandling(e -> e
                    .authenticationEntryPoint((request, response, ex) ->
                            response.sendError(HttpStatus.UNAUTHORIZED.value()))
                    .accessDeniedHandler((request, response, ex) ->
                            response.sendError(HttpStatus.FORBIDDEN.value())))
            .authorizeHttpRequests(a -> a
                    .requestMatchers("/actuator/health", "/api/auth/**").permitAll()
                    .anyRequest().authenticated())
            .addFilterBefore(jwtFilter, UsernamePasswordAuthenticationFilter.class)
            .build();
    }
}
\end{lstlisting}

清单第一行关闭了 CSRF 防护。CSRF（跨站请求伪造）指别的网站诱导用户的浏览器向本平台发请求，浏览器会自动附上本平台的 Cookie，服务端误以为是用户本人在操作；防护办法是要求请求带一个别的网站拿不到的令牌。本例的 API 只从 Authorization 头读取 Bearer 令牌，别的网站的页面拿不到这个头，攻击不成立，所以可以关闭；改用 Cookie 保存凭据时要重新打开它，见本专题“CSRF、CORS 与浏览器存储”部分。清单\ref{lst:online-ch05-jwt-filter}的过滤器只做一件事：把\texttt{Bearer}头里的令牌解析成认证对象放进\texttt{SecurityContext}，解析失败就清空上下文继续放行，由后面的授权规则返回 401，过滤器自己不写响应体。清单开头的\texttt{WaterSystemPermission}枚举是扩展写法：权限按“能做什么”命名，全部写在一个枚举里，名称和取值就不会在各处漂移。

\begin{lstlisting}[label={lst:online-ch05-jwt-filter},language=Java,caption={JWT认证过滤器}]
import jakarta.servlet.*;
import jakarta.servlet.http.*;
import org.springframework.security.authentication.UsernamePasswordAuthenticationToken;
import org.springframework.security.core.authority.SimpleGrantedAuthority;
import org.springframework.security.core.context.SecurityContextHolder;
import org.springframework.security.web.authentication.WebAuthenticationDetailsSource;
import org.springframework.web.filter.OncePerRequestFilter;
import java.io.IOException;
import java.util.stream.Collectors;

enum WaterSystemPermission {
    STATION_READ, READING_WRITE, ALERT_ACKNOWLEDGE
}

class JwtAuthenticationFilter extends OncePerRequestFilter {
    private final JwtService jwtService;
    private final TokenRevocationService revocations;

    JwtAuthenticationFilter(JwtService jwtService,
                            TokenRevocationService revocations) {
        this.jwtService = jwtService;
        this.revocations = revocations;
    }

    @Override
    protected void doFilterInternal(HttpServletRequest request,
                                    HttpServletResponse response,
                                    FilterChain chain)
            throws ServletException, IOException {
        String header = request.getHeader("Authorization");
        if (header != null && header.startsWith("Bearer ")) {
            String token = header.substring("Bearer ".length());
            try {
                Claims claims = jwtService.parseAccess(token);
                String jti = claims.getId();
                if (jti != null && !revocations.isRevoked(jti)) {
                    var authorities = jwtService.permissions(claims).stream()
                            .map(SimpleGrantedAuthority::new)
                            .collect(Collectors.toUnmodifiableSet());
                    var authentication = new UsernamePasswordAuthenticationToken(
                            claims.getSubject(), null, authorities);
                    authentication.setDetails(new WebAuthenticationDetailsSource()
                            .buildDetails(request));
                    SecurityContextHolder.getContext().setAuthentication(authentication);
                }
            } catch (JwtException ex) {
                SecurityContextHolder.clearContext();
            }
        }
        chain.doFilter(request, response);
    }
}
\end{lstlisting}

签名密钥、签发者和令牌有效期这一组值随环境变化，按5.3.2节的办法放在\texttt{application.yml}的\texttt{security.jwt}下，由环境变量提供。清单\ref{lst:ch05-config-validation}用一个\texttt{record}把它们集中起来：\texttt{@ConfigurationProperties}按前缀把配置项绑定到同名字段，\texttt{@Validated}让字段上的校验注解在应用启动时执行。密钥为空、有效期写成 0，应用就拒绝启动并指出是哪一项，不必等到第一个登录请求才暴露。这个\texttt{record}要经过注册才会成为 Bean，清单\ref{lst:online-ch05-security-chain}类上的\texttt{@EnableConfigurationProperties(JwtProperties.class)}做的就是这件事。

\begin{lstlisting}[label={lst:ch05-config-validation},language=Java,caption={JwtProperties：集中绑定 JWT 配置并在启动时校验}]
@ConfigurationProperties(prefix = "security.jwt")
@Validated
public record JwtProperties(
        @NotBlank String secret,
        @NotBlank String issuer,
        @Min(60) @Max(86_400) long accessSeconds,
        @Min(300) @Max(2_592_000) long refreshSeconds) {}
\end{lstlisting}

JWT 密钥是一串足够长的 Base64 值。本章锁定 jjwt 0.11.x：\texttt{parseClaimsJws}在返回 Claims 之前已经完成签名验证，签名不对就抛异常；先手工解码 Base64 再去“检查”字段的写法跳过了签名验证，等于相信了任何人写的令牌。清单\ref{lst:online-ch05-jwt-parse}给出正确的调用顺序，签名验证发生在读取任何声明之前。令牌分两种：访问令牌（\texttt{type=access}）短期有效，每个请求都带；刷新令牌（\texttt{type=refresh}）寿命长得多，只在访问令牌到期时发给刷新端点换一对新令牌，用户不必重新登录。刷新令牌本身到期就拒绝换发，用户重新登录。\texttt{parseRefresh}把刷新令牌交给下面的校验器，它检查签名、类型、签发者、主体、jti 和有效期。

\begin{lstlisting}[label={lst:online-ch05-jwt-parse},language=Java,caption={jjwt 0.11.x 签名验证与访问令牌解析}]
import io.jsonwebtoken.*;
import io.jsonwebtoken.security.Keys;
import io.jsonwebtoken.io.Decoders;
import javax.crypto.SecretKey;
import java.util.List;
import java.time.Clock;
import org.springframework.stereotype.Service;
import edu.example.lesson56.RefreshTokenVerifier;
import org.springframework.context.annotation.Bean;
import org.springframework.context.annotation.Configuration;

// JwtProperties 的定义见上一个清单

@Configuration
class JwtClockConfiguration {
    @Bean
    Clock jwtClock() { return Clock.systemUTC(); }
}

@Service
class JwtService {
    private final SecretKey key;
    private final JwtProperties properties;
    private final Clock clock;

    JwtService(JwtProperties properties, Clock clock) {
        this.properties = properties;
        this.clock = clock;
        this.key = Keys.hmacShaKeyFor(
                Decoders.BASE64.decode(properties.secret()));
    }

    Claims parseAccess(String token) {
        Claims claims = Jwts.parserBuilder().setSigningKey(key)
                .requireIssuer(properties.issuer()).build()
                .parseClaimsJws(token).getBody();
        if (!"access".equals(claims.get("type", String.class)))
            throw new JwtException("token type is not access");
        return claims;
    }

    List<String> permissions(Claims claims) {
        return claims.get("permissions", List.class);
    }

    /** 刷新令牌的签名、声明和有效期统一交给完整校验器。 */
    Claims parseRefresh(String token) {
        return new RefreshTokenVerifier(key, properties.issuer(),
                clock).parse(token);
    }
}
\end{lstlisting}

清单\ref{lst:ch05-refresh-verifier}是这个校验器，文件在配套后端的\texttt{edu.example.lesson56}包里。在\texttt{backend}目录运行\texttt{mvn -Dtest=RefreshTokenVerifierTest test}，可以看到有效令牌通过、到期令牌和篡改过的令牌被拒绝；测试注入固定时钟，到期边界才能反复验证。它只做校验，不是 HTTP 端点；把它接到刷新端点上是清单\ref{lst:ch05-jwt-refresh}的事。

\begin{lstlisting}[label={lst:ch05-refresh-verifier},language=Java,breaklines=true,caption={拒绝到期刷新令牌的完整校验器}]
package edu.example.lesson56;

import io.jsonwebtoken.Claims;
import io.jsonwebtoken.JwtException;
import io.jsonwebtoken.Jwts;
import java.time.Clock;
import java.time.Instant;
import java.util.Date;
import java.util.Objects;
import javax.crypto.SecretKey;

/** jjwt 0.11.x：只验证刷新令牌，轮换与撤销记录由调用方在事务中处理。 */
public final class RefreshTokenVerifier {
    private final SecretKey key;
    private final String issuer;
    private final Clock clock;

    public RefreshTokenVerifier(SecretKey key, String issuer, Clock clock) {
        this.key = Objects.requireNonNull(key);
        if (issuer == null || issuer.isBlank()) throw new IllegalArgumentException("issuer is required");
        this.issuer = issuer;
        this.clock = Objects.requireNonNull(clock);
    }

    public Claims parse(String token) {
        Instant now = clock.instant();
        Claims claims = Jwts.parserBuilder().setSigningKey(key)
                .requireIssuer(issuer).require("type", "refresh")
                .setClock(() -> Date.from(now)).build()
                .parseClaimsJws(token).getBody();
        Date expiry = claims.getExpiration();
        if (expiry == null || !expiry.toInstant().isAfter(now))
            throw new JwtException("refresh token must have a future expiration");
        if (claims.getSubject() == null || claims.getSubject().isBlank())
            throw new JwtException("refresh token subject is required");
        if (claims.getId() == null || claims.getId().isBlank())
            throw new JwtException("refresh token jti is required");
        return claims;
    }
}
\end{lstlisting}

\begin{tipbox}
版本迁移提示：jjwt 0.12.x 将解析器接口拆为\texttt{Jwts.parser().verifyWith(key).build()}与\texttt{parseSignedClaims(token).getPayload()}两步，对应本章 0.11.x 的\texttt{parserBuilder().setSigningKey(key).build()}和\texttt{parseClaimsJws(token).getBody()}。迁移时逐项核对签名算法、issuer、type、过期处理和异常类型，不能只做方法名替换。
\end{tipbox}

认证链路可以拆成凭据校验、令牌签发、请求认证、权限决策、撤销与轮换五个环节。图\ref{fig:online-ch05-auth-sequence}画出访问令牌和刷新令牌各自走的路：访问令牌携带权限进入过滤器；刷新令牌只发送到刷新端点，服务端验证类型和撤销状态后签发新的令牌对。刷新令牌不能当普通 API 凭据用，完整令牌也不出现在浏览器日志、URL 或错误消息里，因为拿到它的人就能冒充用户。


令牌里只放授权用得到的几项声明。\texttt{sub}标识用户或服务主体，\texttt{iss}是签发者，\texttt{aud}是使用方，\texttt{iat}和\texttt{exp}是签发和到期时间，\texttt{type}区分 access 与 refresh，权限集合只带授权决策需要的稳定代码。\texttt{jti}是每个令牌唯一的编号：撤销某一个令牌、在审计里指认某一个令牌，靠的都是它。身份证号、手机号、设备密钥和完整业务对象不放进 JWT，因为 JWT 只是签了名，内容任何人都能解开看，它不是加密容器。

\begin{lstlisting}[label={lst:online-ch05-jwt-issue},language=Java,caption={jjwt 0.11.x 签发访问与刷新令牌}]
import io.jsonwebtoken.Jwts;
import java.time.Instant;
import java.util.Date;
import java.util.List;
import java.util.UUID;

record TokenPair(String accessToken, String refreshToken,
                 Instant accessExpiresAt, Instant refreshExpiresAt) {}

class TokenIssuer {
    private final JwtProperties properties;
    private final SecretKey key;

    TokenIssuer(JwtProperties properties, SecretKey key) {
        this.properties = properties;
        this.key = key;
    }

    TokenPair issue(String subject, List<String> permissions, Instant now) {
        Instant accessExp = now.plusSeconds(properties.accessSeconds());
        Instant refreshExp = now.plusSeconds(properties.refreshSeconds());
        String access = build(subject, permissions, "access", now, accessExp);
        String refresh = build(subject, List.of(), "refresh", now, refreshExp);
        return new TokenPair(access, refresh, accessExp, refreshExp);
    }

    private String build(String subject, List<String> permissions, String type,
                         Instant issuedAt, Instant expiresAt) {
        return Jwts.builder().setId(UUID.randomUUID().toString())
                .setSubject(subject).setIssuer(properties.issuer())
                .setIssuedAt(Date.from(issuedAt)).setExpiration(Date.from(expiresAt))
                .claim("type", type).claim("permissions", permissions)
                .signWith(key).compact();
    }
}
\end{lstlisting}

清单\ref{lst:online-ch05-jwt-issue}给访问令牌和刷新令牌各生成一个 jti；刷新令牌不带权限，权限变更后旧权限只在访问令牌的剩余寿命内残留。签发时间由注入的时钟提供，测试可以固定\texttt{Instant}验证到期边界。生产环境里设备时钟与服务器可能差几秒到几分钟，允许多大偏差写在安全配置里；更换密钥时保留一小段新旧密钥都能验签的窗口，正在处理的请求才不会突然全部失败。清单\ref{lst:online-ch05-login}是登录端点：密码比对走\texttt{PasswordEncoder}，成功后一次签发访问令牌与刷新令牌。配套工程的\texttt{AuthController}是它的单令牌版本，应答只有访问令牌和权限列表。

\begin{lstlisting}[label={lst:online-ch05-login},language=Java,caption={登录端点、密码校验与令牌签发}]
import org.springframework.security.crypto.password.PasswordEncoder;
import org.springframework.web.bind.annotation.*;
import jakarta.validation.Valid;

record LoginRequest(String username, String password) {}
record LoginResponse(String accessToken, String refreshToken,
                     Instant accessExpiresAt, Instant refreshExpiresAt) {}

@RestController
@RequestMapping("/api/auth")
class LoginController {
    private final UserAccountService accounts;
    private final PasswordEncoder passwordEncoder;
    private final TokenIssuer issuer;
    private final Clock clock;

    LoginController(UserAccountService accounts, PasswordEncoder passwordEncoder,
                    TokenIssuer issuer, Clock clock) {
        this.accounts = accounts;
        this.passwordEncoder = passwordEncoder;
        this.issuer = issuer;
        this.clock = clock;
    }

    @PostMapping("/login")
    LoginResponse login(@Valid @RequestBody LoginRequest request) {
        UserAccount account = accounts.findEnabled(request.username())
                .orElseThrow(() -> new BadCredentialsException("用户名或密码错误"));
        if (!passwordEncoder.matches(request.password(), account.passwordHash()))
            throw new BadCredentialsException("用户名或密码错误");
        TokenPair pair = issuer.issue(account.id(), account.permissions(),
                Instant.now(clock));
        return new LoginResponse(pair.accessToken(), pair.refreshToken(),
                pair.accessExpiresAt(), pair.refreshExpiresAt());
    }
}
\end{lstlisting}

数据库里只保存密码的哈希，比对由\texttt{PasswordEncoder}完成。登录失败统一回答“用户名或密码错误”，不区分是用户名不存在还是密码错了：区分了，攻击者就能逐个试出哪些用户名存在。登录端点还要限制调用频率、记录失败次数和审计事件，本专题“密码、账户与登录防护”部分展开；密码重置、二次认证和设备绑定由独立的身份系统负责，不是往 JWT 里多塞几个字段能解决的。刷新端点的检查比登录更容易写漏，清单\ref{lst:ch05-jwt-refresh}把三项必查列全：签名、\texttt{type}是否为\texttt{refresh}、旧 jti 是否已在轮换中作废。

\begin{lstlisting}[label={lst:ch05-jwt-refresh},language=Java,caption={刷新令牌的签名、类型和轮换检查}]
import io.jsonwebtoken.*;
import java.time.Instant;
import java.time.Clock;
import org.springframework.security.authentication.BadCredentialsException;

// 实现须以事务中的条件更新消费未过期且主体匹配的会话，返回旧会话；
// 同一 jti 的并发调用至多一个成功。登录时也须登记新 refresh 会话。
record RefreshSession(String subject, java.util.List<String> permissions) {}
interface RefreshSessionRepository {
    java.util.Optional<RefreshSession> consumeActive(
        String jti, String subject, Instant now);
    void saveActive(String jti, RefreshSession session, Instant expiresAt);
}

@org.springframework.stereotype.Service
class RefreshService {
    private final JwtService jwtService;
    private final TokenIssuer issuer;
    private final TokenRevocationService revocations;
    private final RefreshSessionRepository sessions;
    private final Clock clock;

    RefreshService(JwtService jwtService, TokenIssuer issuer,
                   TokenRevocationService revocations,
                   RefreshSessionRepository sessions, Clock clock) {
        this.jwtService = jwtService;
        this.issuer = issuer;
        this.revocations = revocations;
        this.sessions = sessions;
        this.clock = clock;
    }

    @org.springframework.transaction.annotation.Transactional
    TokenPair refresh(String token) {
        Claims claims = jwtService.parseRefresh(token);
        String jti = claims.getId();
        if (jti == null || revocations.isRevoked(jti))
            throw new BadCredentialsException("刷新会话已撤销");
        Instant now = Instant.now(clock);
        RefreshSession session = sessions.consumeActive(jti, claims.getSubject(), now)
                .orElseThrow(() -> new BadCredentialsException("刷新会话不存在"));
        revocations.revoke(jti, claims.getExpiration().toInstant());
        TokenPair next = issuer.issue(session.subject(), session.permissions(), now);
        Claims nextClaims = jwtService.parseRefresh(next.refreshToken());
        sessions.saveActive(nextClaims.getId(), session,
                nextClaims.getExpiration().toInstant());
        return next;
    }
}
\end{lstlisting}

JWT 的声明字段、签名与校验语义以 RFC 7519 为准\cite{rfc7519}。刷新采用轮换：旧刷新令牌的 jti 用过一次就作废，同时发一个新的。同一个旧令牌第二次出现，说明它被复制过，这叫重放，此时可以把这个用户的整个会话族一起撤销。到期的刷新令牌直接拒绝。\texttt{consumeActive}用一条带条件的 UPDATE 原子地把旧会话标为已用，并在同一个数据库事务里登记新会话，两个并发的刷新请求最多一个成功；“先查询再写 Redis”做不到这一点，两个请求可能都查到“未用过”。这个仓储按5.4.6节的条件更新实现；Redis 里的撤销记录只是缓存，会话事实在数据库里。

\begin{lstlisting}[label={lst:ch05-revocation},language=Java,caption={基于 jti 与版本号的令牌撤销}]
import java.time.Duration;
import java.time.Instant;
import org.springframework.data.redis.core.StringRedisTemplate;

interface TokenRevocationService {
    boolean isRevoked(String jti);
    void revoke(String jti, Instant expiresAt);
}

@Service
class RedisTokenRevocationService implements TokenRevocationService {
    private final StringRedisTemplate redis;

    RedisTokenRevocationService(StringRedisTemplate redis) {
        this.redis = redis;
    }

    public boolean isRevoked(String jti) {
        return Boolean.TRUE.equals(redis.hasKey("jwt:revoked:" + jti));
    }

    public void revoke(String jti, Instant expiresAt) {
        Duration remaining = Duration.between(Instant.now(), expiresAt);
        if (!remaining.isNegative() && !remaining.isZero())
            redis.opsForValue().set("jwt:revoked:" + jti, "1", remaining);
    }
}
\end{lstlisting}

JWT 签出以后服务端不保存它，到期之前一直有效。用户注销或改密后要让令牌提前失效，只能把它的 jti 记进一份“已撤销”名单，过滤器每次都查一下，这就是令牌撤销。清单\ref{lst:ch05-revocation}把名单放在 Redis 里，每条记录的 TTL 等于原令牌的剩余寿命，令牌到期后记录自动消失，名单不会无限增长。另一种做法是给每个用户记一个令牌版本号：注销或改密时版本加一，过滤器比较令牌里的版本和账户当前版本。jti 名单适合撤销单个设备，版本号适合撤销用户的全部会话，两者可以并用。Redis 连不上时按本专题“撤销、重放与故障降级”部分的策略处理：写操作拒绝，只读请求可以放行；不能把“查不到”当作“没有撤销”。

过滤器链解决“是谁”，方法级注解解决“能做什么”。清单\ref{lst:online-ch05-method-security}用\texttt{@PreAuthorize}把权限要求写在控制器方法上，权限名就是配套\texttt{SecurityConfig}里三个教学账号的 DUTY、ANALYST、OPS。清单后半是跨域配置。CORS（跨源资源共享）是浏览器的一条规则：页面向别的源（协议、域名或端口有一项不同）发请求时，浏览器先问对方服务器是否允许这个源读取应答，服务器用响应头回答；开发时 Vite 在 5173 端口、后端在 8080 端口，就是这种情况。跨域策略集中写在一处，不散落在各个控制器里。

\begin{lstlisting}[label={lst:online-ch05-method-security},language=Java,caption={方法级权限与跨域策略}]
import org.springframework.security.access.prepost.PreAuthorize;
import org.springframework.web.bind.annotation.*;

// AssetController.java（配套工程）：查询接口对三种角色开放
    @GetMapping("/{assetId}/readings/latest")
    @PreAuthorize("hasAnyAuthority('DUTY','ANALYST','OPS')")
    public ResponseEntity<ReadingDto> latest(@PathVariable String assetId) { /* 见 5.4.2 节 */ }

// ReadingWriteController.java（5.5.1 节）：补录只对专业分析员开放
    @PostMapping
    @PreAuthorize("hasAuthority('ANALYST')")
    public ResponseEntity<AssetController.ReadingDto> create(/* 见 5.5.1 节 */) { /* …… */ }

@Configuration
class CorsConfig {
    @Bean
    CorsConfigurationSource corsConfigurationSource() {
        CorsConfiguration c = new CorsConfiguration();
        c.setAllowedOrigins(List.of("https://water.example.edu"));
        c.setAllowedMethods(List.of("GET", "POST", "PUT", "DELETE"));
        c.setAllowedHeaders(List.of("Authorization", "Content-Type"));
        UrlBasedCorsConfigurationSource source = new UrlBasedCorsConfigurationSource();
        source.registerCorsConfiguration("/api/**", c);
        return source;
    }
}
\end{lstlisting}

方法级权限是前端按钮隐藏之外的第二道边界：按钮藏起来只是看不见，直接发请求照样能到后端。CORS 只约束浏览器能不能读跨源应答，和身份认证是两回事；允许的来源、方法和请求头写成明确的白名单，带凭据的请求不能允许任意来源。401 表示没有认证，403 表示认证了但没有权限，前者页面要跳登录页，后者不跳，所以两者不能混用。日志记权限代码和追踪 ID，不回显完整 JWT。

配套工程只到这里：三个账号、三种权限、一种令牌，用\texttt{smoke.sh}或第4章的页面就能看到 401 和 403 的区别。认证与授权的测试写法见5.9.2节，\texttt{MockMvc}配合\texttt{@WithMockUser}模拟登录身份，覆盖合法令牌、无令牌和权限不足三种情形。后面的扩展写法各自还要补测试：签名错误、issuer 错误、type 错误、过期令牌、撤销 jti，以及旧 refresh jti 第二次使用被拒绝；测试用随机生成的密钥和固定时钟，不把真实密钥写进源码。

\subsection{威胁模型与认证边界}

安全设计先列出要保护的东西和攻击者可能走的路，这份清单叫威胁模型。案例水库平台要保护值班员身份、测站写入权限、预警处置记录、设备上报密钥和审计日志；攻击者可能伪造 Bearer 令牌、重放偷来的刷新令牌、利用错误信息的差异试探账号、用跨源页面诱导浏览器发请求，或者在日志和备份里翻找长期有效的密钥。JWT 只回答“令牌是谁签的、内容有没有被改”；令牌被复制后的重放、传输加密、账户停用和数据库权限，都要靠别的机制。

访问令牌的有效期以泄露后可接受的暴露时间为上限，配套工程取 1800 秒；刷新令牌的寿命由设备风险、撤销能力和换班流程决定。有效期长，泄露后能被利用的时间就长；有效期短，客户端刷新得频繁，网络不稳时更容易被登出。取值时把网关超时、Redis 可用性、时钟偏差和移动网络重连一起考虑，写进配置和安全评审记录，不写成代码里的常数。

\subsection{密钥生命周期与配置管理}

签名用的 HMAC 密钥由密钥管理器或部署平台的机密变量提供，配套工程用环境变量\texttt{JWT\_SECRET}或\texttt{JWT\_SECRET\_FILE}注入（8.6.1节），代码里那个默认值只供课堂演示。应用启动时检查密钥能否 Base64 解码、长度是否够算法要求、issuer 是否为空，有一项不对就拒绝启动，比在第一个登录请求上失败好定位。日志只输出密钥版本和摘要，不输出密钥本身、完整 JWT 或配置对象的\texttt{toString}。

更换密钥分四步：生成新密钥并登记版本；验签同时接受新旧两把（按令牌头部的\texttt{kid}选择），签发只用新的；观察认证失败率，撤销异常版本；等最长的刷新令牌寿命加时钟容差过去以后，删除旧密钥。令牌没有\texttt{kid}时，验签只能按时间窗口逐个尝试候选密钥，下一次签发就把版本标识补上。

\subsection{密码、账户与登录防护}

密码不存明文，存的是 Spring Security 的\texttt{PasswordEncoder}算出的哈希。配套工程通过\texttt{PasswordEncoderFactories}类的\texttt{createDelegatingPasswordEncoder()}方法取得编码器，默认算法是 BCrypt，哈希值前面带着算法标识，将来换算法时旧记录仍能比对。数据库只保存哈希、算法版本、修改时间和失败计数，不保存明文，也不保存可以解回明文的密文。登录成功后失败计数清零，并记录设备和追踪 ID；连续失败后逐次延长等待时间或临时锁定。运维员重置密码时把账户的令牌版本加一，这个用户已签出的访问令牌和刷新令牌就全部失效。

认证失败的对外表现要一致。用户名不存在、密码错误、账户停用和二次认证失败都返回同样的 401 描述，具体原因只写进审计；资源是否存在的检查放在权限判断之后，否则没有权限的用户可以从 404 和 403 的差别猜出哪些测站编号存在。审计事件记录主体、客户端、时间、结果码和\texttt{jti}，不记录密码、Authorization 头或完整刷新令牌。

\subsection{CSRF、CORS 与浏览器存储}

本章过滤器链关闭 CSRF，前提是 API 只从 Authorization 头读令牌，浏览器不会自动附带它。改用 HttpOnly Cookie 保存凭据以后，浏览器对每个请求都会自动带上 Cookie，别的网站的页面也能借用户的浏览器发出带 Cookie 的请求，这时就要打开 CSRF 令牌或采用双重提交方案。HttpOnly 只是让脚本读不到 Cookie，挡不住跨站请求本身。Cookie 的 SameSite、Secure、过期时间和域路径与部署拓扑一起定。

CORS 预检是浏览器在真正发请求之前先发一个\texttt{OPTIONS}请求问服务器“允许这个源吗”，服务器答应了，浏览器才发真正的请求；真正的请求照样要经过认证和授权。生产环境把来源写成白名单，明确允许的方法、请求头、暴露头和缓存时长。预检失败时正确的做法是把来源加进白名单，把允许来源改成\texttt{*}等于关掉这道检查。非浏览器的设备、Kafka 消费者和内部任务不经过浏览器，CORS 管不到它们，它们靠服务身份和网络策略。

令牌放在浏览器的哪里，各有代价。放在内存里最安全，但页面一刷新就没了；\texttt{sessionStorage}随标签页存在，同源脚本能读；\texttt{localStorage}能持久登录，页面一旦被注入脚本（XSS），令牌就长期暴露。要求更高的系统把刷新令牌放进受保护的 Cookie，访问令牌放内存，再配合内容安全策略、依赖审计和严格的 DOM 写入规则。不管选哪种，网络日志、错误监控和浏览器历史里都不该有令牌。

\subsection{权限建模与最小授权}

配套工程把角色名直接当权限用：DUTY、ANALYST、OPS 各对应一个账号。系统变大以后，权限按“能做什么”命名比按角色命名稳定：\texttt{STATION\_READ}允许读测站，\texttt{READING\_WRITE}允许写观测，\texttt{ALERT\_ACKNOWLEDGE}允许确认预警；值班员、专业分析员、审批人和运维员是角色，角色拥有哪些权限由服务端配置决定，调整时改配置，不改代码。资源范围要单独检查：有写观测的权限，不等于能改所有工程的观测或归档记录。

方法级注解适合粗粒度的“能不能做这类事”，测站归属、时间窗口和状态迁移这些细粒度判断放在应用服务里。授权失败要在任何写入之前发生，先保存实体再检查权限，拒绝了也留下了副作用；批量接口逐条检查资源范围，一条请求不能顺带改到别的测站。增加权限向后兼容，删除权限或改变含义要发布新版本，审计里记下当时用的授权策略版本。

\subsection{撤销、重放与故障降级}

撤销名单在 Redis 里，Redis 连不上时怎么办要事先定好。写入、改密和预警处置这类高风险操作拒绝并返回 503；读取公开状态可以用短时间的缓存，应答里标明缓存时间。名单记录带 TTL，至少覆盖令牌剩余寿命；用户版本号在数据库事务里更新，缓存过期后自然一致。撤销接口本身也要认证、幂等和审计，否则任何人提交一个 jti 就能把别人踢下线。

重放检测不只用于刷新令牌，一次性登录链接、密码重置链接和设备注册码都是同一类问题：服务端保存随机值的哈希、主体、用途、首次使用时间和过期时间，第一次使用时原子地标记已用，第二次使用返回统一的错误。HTTP 重试和消息重投带请求 ID，按 ID 判断重复；按到达时间猜不可靠。同一个刷新会话族出现异常重放时，撤销这个用户的全部会话并通知运维员复核。

\subsection{认证链路的可观测性与演练}

认证链路要观察的指标有 401/403 的比率、签名失败次数、issuer 或 type 不匹配次数、刷新成功率、撤销查询延迟、Redis 错误、密钥版本分布和账户锁定次数。令牌解析失败的原因分类只进指标，不返回给客户端：签名错、过期、issuer 不符、type 错和已撤销分别计数，应答统一是 401；这样攻击者得不到试错反馈，运维员却能从指标看出是密钥配错了、时钟漂了还是有人在重放。指标标签不用完整用户名、令牌或原始 URL，否则监控系统自己就成了泄露点。日志用结构化字段记\texttt{traceId}、主体哈希、jti 哈希、权限结果和策略版本；审计日志单独存放，另设保留期和访问权限。

演练要覆盖密钥配错、Redis 故障、旧密钥退役过早、时钟漂移、刷新令牌重放和网关改写错误码。每次记录发现用了多久、拒绝了多少请求、怎样恢复、哪些请求类型没覆盖到；恢复后检查旧令牌是否还能访问、撤销名单有没有异常增长、缓存里有没有过期的权限。接口契约里写明 Authorization 头的大小上限、令牌格式、401/403/429/503 的应答结构和重试建议。限流分三层：网关做基础限流，应用服务做账户和 jti 级别的风控，数据库靠唯一约束保证审计不重复，每层都不假定别的层已经做了。

\subsection{令牌失效与会话迁移}

退出登录要同时处理浏览器和服务端两边。客户端清掉内存里的访问令牌，调用注销接口撤销当前刷新令牌的 jti，之后收到 401 就不再自动重试，否则失败请求会循环下去。服务端的注销接口先核对这个刷新会话属于当前用户和设备，再写撤销记录；运维员强制某人下线，则把账户版本加一并记下原因。一个用户在多台设备上登录时，注销一台不影响其他台，除非用户明确选择“退出全部会话”。

权限变更不会立刻反映到已签出的令牌上。用户从“可写入”改成“只读”以后，手里的访问令牌在剩余寿命内仍带着旧权限。高风险操作有三种办法缩短这个窗口：访问令牌有效期更短，过滤器实时比较账户版本，或者敏感接口在服务层再查一次账户状态。普通读取可以直接相信令牌里的声明，预警确认、规则发布和测站配置修改在服务层用最新权限判断。

刷新端点限制调用频率和来源。客户端只在访问令牌临近到期或收到明确的 401 时刷新一次，多个并发请求用一把互斥锁合并成一次刷新，否则它们会同时去消费同一个旧 jti，只有一个能成功，其余全部被登出。轮换成功后客户端只保留最新的一对令牌；轮换失败就清掉本地凭据回到登录页，被撤销的旧令牌不能继续用。

审计数据分两类：登录成功、登录失败、刷新、撤销、密钥切换是认证事件；创建观测、确认预警、修改配置是业务操作。两类共用 traceId 和主体，保留期、可读角色和脱敏规则各自设定。审计表以事件 ID 为唯一约束，写审计失败时按5.4.5节的事务边界决定是回滚主操作、写入发件箱还是转人工补偿，不能悄悄丢掉。

多个服务是否共用一把 JWT 签名密钥要明确决定。共用一把验签方便，泄露影响也大，任何一个服务都能伪造其他服务信任的令牌；改用非对称算法集中签发，各服务只拿验签用的公钥，签发服务保护好私钥。不管用 HMAC 还是非对称算法，issuer、audience、允许的算法列表和密钥版本都写进契约，绝不按令牌头部里写的算法名去信任一个没见过的算法。

写这部分代码时容易犯的错，集中在几处：手工 Base64 解码后直接相信 Claims；用没有验签的解析方法；把 refresh 当 access 用；把密钥放进前端变量或提交到 Git；异常信息里回显完整令牌；用放宽 CORS 代替权限；在过滤器里做数据库写事务。配套工程能跑通，只说明这条链的骨架是对的；上面这几条一处没守住，认证就形同虚设。

\section{前后端认证联调的深入检查}\label{app:ext-fe-auth}\label{online:frontend-auth}

\paragraph{对应正文}第4章4.5.7—4.5.8节与第5章5.6节。

本节接第4章4.5.7、4.5.8节，需要第5章5.6节的认证与授权。

\paragraph{先画状态机}认证相关的页面状态至少有六种：匿名、已登录、令牌即将过期、令牌已过期、权限不足和服务不可达。每种状态都要规定界面文案、是否保留原路径、是否允许重试，以及日志里应出现的事件。例如匿名用户访问测点详情，守卫保存\texttt{redirect=/assets/DAM-A-PZ-07}并跳转登录，登录成功后只恢复一次导航；后端返回 403 时停留在业务页面并显示权限说明，不再跳到登录页。令牌过期后页面在登录页和业务页之间反复跳转，通常是把这六种状态折叠成了“有令牌”和“没令牌”两种。

\paragraph{登出与多标签页}登出除了删除浏览器里的令牌，还要调用后端的撤销接口，或者让短期令牌自然过期；多个标签页通过\texttt{storage}事件同步退出状态。刷新令牌如果放在 HttpOnly Cookie 里，前端读不到它，只能根据刷新接口的结果更新内存中的访问令牌。

\paragraph{字段升级的节奏}前后端对字段的兼容按“先增加、后切换、再清理”进行。后端新增\texttt{qualityFlag}时，一段时间内同时返回新旧字段；前端先读新字段、缺失时回退旧字段，并统计回退次数；确认所有客户端升级后再删除旧字段。删除之前检查缓存、导出任务和消息消费者。日期时间统一用带时区的 ISO 8601 字符串，数值单位写进字段说明，枚举值变化由版本化接口或兼容映射处理。组件不直接拼接 URL、读取令牌或解释状态码，这些集中在请求模块里；路径前缀或分页字段调整时，只改这一处。

\paragraph{跨域与预检}开发时浏览器向 Vite 发同源请求，再由代理转发。生产环境如果前端与接口不同源，允许的来源、方法、请求头和是否允许凭据由网关与后端共同约定。带\texttt{Authorization}头的请求会先发一次\texttt{OPTIONS}预检，服务端要正确应答预检，但预检通过不等于真正的请求可以跳过认证。\texttt{Access-Control-Allow-Origin: *}不能与带凭据的请求同时使用，允许来源应是配置好的白名单，不能把请求里的\texttt{Origin}原样返回。遇到“浏览器报跨域、curl 正常”时，分别检查预检响应、实际响应和网关是否丢掉了认证头。

\paragraph{令牌放在哪里}\texttt{sessionStorage}便于教学演示，标签页之间互相隔离，但同源脚本可以读取它；HttpOnly Cookie 让脚本读不到令牌，代价是服务端要设置 SameSite、Secure 并防护 CSRF。无论哪种方式，访问令牌都不能出现在 URL、页面标题或埋点参数里；浏览器控制台只输出状态码、错误码和追踪 ID，服务端日志对令牌做哈希或只记录 jti；前端构建产物中不能出现 JWT 签名密钥、数据库连接串和内部网段。

\paragraph{认证的性能}一次受保护请求的耗时分为令牌解析、撤销查询和业务查询三段。每次请求都访问远程撤销存储会增加延迟，完全不查又无法及时吊销泄露的令牌；折中做法是访问令牌设为短期，把 jti 撤销集合缓存在 Redis，过期时间与令牌的 exp 对齐。角色到权限的展开在服务端完成，前端据此隐藏不该出现的按钮，但“按钮不可见”不能代替后端授权。批量读取测点时，后端限制单次数量和时间窗口，避免一个合法令牌被用来放大高负载接口。

\paragraph{可重放的联调记录}固定测试数据和时间窗口：每个测点至少准备一条 valid、一条 suspect 和一条 missing 读数，分别验证正常展示、可信度降低和不参与判断；告警处置单准备“待确认—已确认—已关闭”三种状态，检查不同角色能执行的动作。测试请求携带固定的\texttt{X-Request-Id}，应用日志保留它以及测点编码、质量码和耗时，数据库审计记录补充操作者和事务结果。定位一次失败请求时，先记下用户动作和路由，再记请求方法、路径、状态码、错误码和追踪 ID，最后核对过滤器日志、业务日志与审计表是否指向同一操作者和测点。联调结束后删除调试开关和临时账号，把前端构建版本、后端提交号、数据库迁移版本、代理目标和测试账号角色一并归档，下次升级照同一份记录重放，就能分清问题来自新代码还是环境变化。

\paragraph{按角色走查}在案例水库的演练中，值班员从测点列表打开详情，专业分析员提交处置意见，审批人确认后才允许关闭事件。三个角色看到的按钮、请求权限和审计结果必须一致；任何一个角色在页面上看到可点击的操作却收到 403，就回到路由元数据、后端方法授权和角色映射逐层核对。


请求并发时的“竞态覆盖”已在 4.5.4 节用清单\ref{lst:ch04-r1-race}处理；自动化测试可以让第一个模拟请求延迟 200 毫秒、第二个只延迟 20 毫秒，再快速调用两次\texttt{showAsset}，断言页面最终文字属于第二个对象、第一个请求收到取消信号。

发布前可以进行一次“故障注入式”验收：把响应状态依次改为 401、403、404、409、429 和 503；把响应延迟设置为 100 毫秒、超时阈值和超时阈值的两倍；将令牌的 exp 调整到过去、将 issuer 改成错误值、删掉 permissions 字段；再让代理返回无效 JSON。每个场景都要记录页面是否可操作、是否产生重复请求、是否保存了原路径、是否有泄露敏感信息和是否生成唯一 traceId。每次注入故障后，页面都应保留可用操作、说明失败原因并提供恢复入口，同时将故障与正常水位、空告警结果明确区分。

定位联调故障时，可先用状态码判断请求结果，再用错误码查找业务原因；按 traceId 检索跨服务时间线，最后核对审计记录中的写入结果。值班员提交的测站编号、发生时间和页面提示，应能对应到这些记录。

\section{从业务约束到物理设计的完整推导}\label{online:db-design}

\paragraph{对应正文}第8章8.3.7节（从业务约束到物理设计的推导）。正文保留业务事实到字段、事务边界和索引推导三段；本节收录外键、枚举与数值约束、时序分区、连续聚合、备份、权限、模式演进、测试与文档的完整论述。

外键是数据库层最直接的业务保证。\texttt{warning.asset\_id}要求预警必须指向仍可识别的工程对象；
\texttt{work\_order.warning\_id}要求工单不能脱离预警凭空出现；模型运行不直接外键到每一条观测，
而是保存输入快照，因为一次运行可能读取聚合结果、多个测点和一组规则参数。若把模型输入硬编码成
若干外键，后续版本的输入范围难以重现。快照中的事件范围、质量过滤条件和查询版本必须写入可审计字段，
并在模型日志中登记哈希摘要。

枚举约束的价值在于让错误尽早暴露。质量码只允许 valid、suspect、missing，预警等级只允许 NONE、
BLUE、YELLOW、ORANGE、RED，工单和模型运行也各自拥有有限状态集合。前端可以提供友好文案，但不能
向数据库写入“正常”“蓝色预警中”等自由文本。新增状态先扩展接口契约、迁移脚本和统计口径，再修改
显示层；如果只改显示层，历史报表会出现同义词，统计无法按状态合并。

数值约束必须与测量物理意义相符。观测值允许为空只在质量码为 missing 时成立；温度、位移、渗压等
监测项的合理范围由测点类型和项目规则提供，通用表只负责保留数值和质量结论，不把某个工程的阈值
硬编码到所有测点。单位字段保留在观测行中，是因为设备更换或校准可能改变单位；单位注册表负责检查
一致性，历史数据不能因为后来修改注册表就失去原来的解释。修订时新增版本并携带原因，审计可以看出
从何值改到何值、谁批准、何时生效。

时序分区的核心是把时间局部性转化为物理局部性。查询最近一小时的测点，不应触碰数年前的块；归档
两年前的数据，不应锁住实时写入块。切块大小需要用代表性写入压测确定：块太小会让规划器管理大量
对象，块太大则使压缩和备份粒度变粗。哈希分区解决的是热点测点写入集中问题，不能代替时间维度；两者
同时使用时，要观察单测点查询是否因为跨分区而退化。生产配置和压测结果必须写入部署文档，不把例题
中的四个分区误当作全书固定参数。

连续聚合适合重复读取的统计结果，但不适合替换证据。预警计算可以读取十五分钟平均值以降低延迟，
但触发处置前要保留对应的原始窗口、样本数和质量分布。聚合刷新存在迟到数据，因此策略需要保留一段
回看窗口；通信链路补传超过窗口时，运维员应运行一次受控重算，并在\texttt{model\_run}或变更日志中
注明时间范围。任何人都不能直接编辑聚合表来“修正报表”，修正应回到原始观测或重新运行聚合策略。

备份策略要同时覆盖结构、数据、权限和扩展。PostGIS 的空间列、TimescaleDB 的超表元数据、连续
聚合策略和角色授权缺一项都可能导致恢复后查询失败。全量备份之外还要保留增量或归档日志，并定期在
隔离环境恢复。恢复后按事件标识查询测点、原始观测、质量结论、预警记录、工单回执和模型输入，
再检查角色授权是否正确、新的时间块能否继续写入；任一项失败，都应记录恢复缺口。

权限设计遵循最小必要原则。采集角色只写观测和发件箱，质量角色可更新质量标记但不能修改原始值，
预警角色可创建和关闭预警，值班员可确认工单并提交回执，分析员可读取脱敏的历史数据，运维员才可
执行分区、压缩和恢复操作。数据库角色与应用角色一一对应，凭据存放在密钥服务中，代码和教材示例不
出现真实口令。审计日志至少保留角色、对象、操作、时间、请求追踪号和结果摘要，便于把数据库行为与
Kafka 消息、HTTP 请求和前端操作对应起来。

模式演进要优先考虑兼容。增加可空列通常可以向后兼容，删除列和改变枚举则必须经过观察期；大表添加
非空约束应先回填并验证，再分阶段切换。实体类和 DTO 不应复制一套互不相同的字段名，数据库字典是
唯一来源，代码生成或手工实现都应在代码审查中逐列核对。若后端需要把\texttt{asset}映射为 Java
实体，应保留\texttt{asset\_id}、\texttt{geometry}和\texttt{metadata}的原语义，并对空间类型
明确转换策略；若 S3 生成示例数据，应按照同一编码规则和质量码集合生成，不能为方便而使用随机中文键。

测试要覆盖约束、查询计划和故障恢复三个层次。单元测试验证质量优先级和分类结果，集成测试验证外键、
唯一键、空间坐标系、超表写入和连续聚合，端到端测试验证从事件接入到工单回执的链路。性能测试使用
与目标测点规模相近的数据分布，分别记录高峰写入、最近查询、空间查询、预警队列和聚合刷新延迟。
故障测试包括消息重复、乱序、数据库短暂不可用、超表块压缩失败和恢复后重放；每个故障都要说明预期
质量码、状态转移和用户可见提示。

最后，数据库文档必须让不在场的同学也能复现。字段字典说明类型、单位、约束和示例值，迁移脚本说明
版本顺序和回滚边界，ER 图说明关系方向，查询练习说明输入参数和结果含义，运维手册说明备份、恢复、
压缩、保留和监控。课程项目验收时，学生提交一份从\texttt{event\_id}出发的追溯报告，附执行计划、
质量统计、预警证据、工单回执和模型运行快照。报告中若出现“数据库里有记录但无法解释来源”，就说明
数据契约尚未闭合；若出现“界面显示正常但质量码为 missing”，就说明展示层没有尊重数据库语义。

实体映射与接口实现应沿用这组数据定义：对象表提供稳定编码，观测表保留时间与质量，预警和工单表
记录评估与处置，模型运行表保存输入快照和版本。实现时逐列核对Java类型、JSON字段和数据库约束，
再用查询练习检查相同输入是否得到一致结果。

在课堂演练中，可以把一次数据异常拆成四个观察窗口。第一窗口查看接入消息，确认事件 ID、测点编码
和发生时间是否完整；第二窗口查看质量服务，确认缺测、单位错误、范围异常、时序倒置和重复事件是否
分别留下原因；第三窗口查看数据库，确认原始观测只追加、预警证据包含规则版本、工单具有责任角色和
时限；第四窗口查看前端，确认“无预警”“蓝色关注”和“未评估”使用不同文案和颜色。四个窗口的记录
必须能够用追踪号互相定位，单独看某一层的截图不能代替链路证据。

数据库容量规划也应从测点数量、采样周期和保留期限推导。估算时先计算每天的事件行数，再加入索引、
JSONB 证据、备份副本和压缩后的增长系数；写入峰值还要考虑设备同时上线、补传和消息重投。估算结果
只用于容量初始值，最终仍要用接近生产分布的回放数据压测。压测不仅看平均吞吐，还要记录第九十五和
第九十九百分位延迟、锁等待、超表块数量、连续聚合刷新滞后以及故障重试后的积压恢复时间。

当项目需要接入新的监测项时，先在数据字典登记名称、单位、精度、采样周期和质量规则，再创建或复用
测点类型；不要直接在 JSONB 中写一个未经登记的键。新增字段若会参与预警、统计或模型输入，必须同步
更新接口契约、索引评估、备份清单和验收查询。这样做虽然增加了前期沟通，却能避免后期出现“报表有一
列、接口没有一列、模型又按另一列计算”的隐性分叉。

数据库契约还应服务于教学复盘。学生可以从事件追溯查询出发，逐个删除或改变查询条件，观察外键、质量
码、状态和索引如何影响结果；也可以故意发送重复事件、未来时间戳和缺少单位的记录，验证系统是否在
正确的位置拒绝或标记。实验报告要写清“哪一层发现问题、哪一层保存证据、哪一层向用户解释”，而不
只是贴出一条 SQL 的执行结果。将观察结果对应到外键、状态约束和查询条件，说明各项约束的作用。

工程交付时还要把数据库契约交给不同角色共同评审：业务人员检查字段是否能表达处置流程，测量人员
检查单位和精度，开发人员检查实体与 DTO 映射，运维人员检查分区、备份和恢复，安全人员检查角色与
审计范围。评审结论写入版本记录，后续修改沿用同一流程。进入实际项目时，
字段类型、关系与约束可作为设计起点；工程阈值、角色授权、分区规模和保留期限应按项目规程、审批结果
与压测数据重新核定。

因此，数据库评审记录应成为发布包的一部分，与代码、配置和数据样例共同归档，并在每次迁移后更新。
归档内容还应包含执行计划、失败重试记录和恢复演练的校验结果，保证后来接手系统的人员能够复现设计
依据并核对变更后的查询结果与性能。这些材料同时构成课程答辩的证据链：学生在报告中引用对应的
迁移版本和查询清单，并注明每条证据的生成时间，评审者即可逐项复核。

数据库设计还要考虑数据进入系统之前的契约校验。网关收到消息时先检查事件标识、测点编码、发生时间、
单位和载荷格式；缺少任一必填项就进入拒收或隔离队列，并保留原始报文摘要。通过校验的消息再进入质量
服务，质量服务依据测点类型读取量程和单位注册表，生成 valid、suspect 或 missing 质量码。这个顺序
可以把“消息格式错误”和“数值虽然可解析但不可信”区分开，后续统计和预警才不会把两类问题混为一谈。

观测写入采用追加式版本模型，会牺牲少量存储空间，却换来清晰的审计边界。设备重传同一事件时，唯一
事件键保证幂等；人工确认量程或修复时间戳时，新增版本并在修订原因中说明依据。查询最新值必须明确
按版本和事件时间排序，报表若要重现过去状态，则读取指定版本快照。这样既能支持现场快速纠错，也不会
让“修正后的数字”覆盖掉审计人员需要的原始事实。

空间数据的精度需要与业务用途匹配。点位用于测点定位和影响区筛选时，使用带高程的 PointZ 能够同时
满足平面地图和剖面分析；如果某个业务需要线或面，应在项目扩展表中登记几何类型、坐标系和精度要求，
不要把线面强行塞进点字段。三维模型中的局部坐标、工程坐标与经纬度之间要保留转换参数，并在模型
运行输入快照中登记版本，避免相同测点在不同软件中出现肉眼难以发现的偏移。

时间数据的准确性也不能只看时间戳格式。设备时钟漂移、网关缓存和消息重投都会造成到达顺序与发生顺序
不一致。数据库同时保存 \texttt{occurred\_at}、\texttt{received\_at} 和事件 ID，质量服务依据当前时间与上一条观测检查
未来时间、倒序和重复；流式消费者按发生时间排序时必须设置允许迟到窗口。窗口外的补传进入人工复核，
而不是静默插入实时曲线。通过把三个时间概念分别保存，系统可以解释“数据什么时候发生、什么时候到达、
什么时候被平台处理”。

预警表的证据字段应能让专业分析员复算当时的判定。证据至少包括归一化指标、权重、阈值、规则版本、
参与计算的数据窗口和质量统计；只保存一个分数会使处置过程无法复核。无预警意味着规则已经成功计算
且没有触发阈值，未评估意味着质量条件不足，蓝色则表示指标出现需要关注的变化。三者在数据库中分别
由等级、evaluable 和 reason 表达，前端不得把空 reason 当作正常状态。

工单记录预警发生后的处置过程。预警可能被多次确认、转派和补充回执，工单需要记录责任角色、
截止时间、动作、结果和完成时间。状态更新必须使用当前状态条件和追踪号，若并发更新失败则返回最新状态
供用户刷新。处置结果应包含测点、时间、操作前后状态和依据，不能只写“已处理”。这些字段支持事后
统计处置时延和重复工单，也能帮助模型评估哪些预警经常被人工驳回。

模型运行记录把算法版本和输入快照固定下来，解决“同样的代码为什么结果不同”的问题。输入快照包括
测点集合、时间范围、质量过滤条件、特征版本和规则参数；结果字段存储摘要，完整输出可存入版本化对象
存储并在表中保存地址和校验和。运行失败时记录错误类别和重试次数，成功后才能被预警或报表引用。模型
升级必须创建新版本并进行对照运行，旧版本结果继续保留，直到审批人确认迁移完成。

最后，数据库指标应反映用户关心的业务体验。除了写入吞吐和查询延迟，还要观察最近数据滞后、合格数据
比例、未评估预警数量、待处理工单时长、连续聚合刷新延迟和恢复演练成功率。监控看板把这些指标按测点、
工程和时间窗分组，运维员可以区分“设备没有上报”“数据到达但质量不合格”“规则未运行”和“工单没有
及时关闭”。不同状态对应不同处置入口：检查设备、复核数据、恢复规则任务或催办工单。

\section{平台形态并存时的隔离、抽查与评审方法}\label{app:ext-ch01-platform-review}\label{online:ch01}

\paragraph{对应正文}第1章1.2节。

\subsection{“问题—尺度—证据”卡片与租户隔离}

形态选型可以用一张“问题—尺度—证据”卡片完成。问题栏写清需要解决的业务动作，尺度栏写清涉及的工程、管理站、区域或流域范围，证据栏写清需要看到的字段、时延、权限和回执。例如“管理站是否能在网络波动时完成配水计划复核”属于区域与工程交界问题，既要有本地缓存和数据龄期标记，也要有计划版本、审批人和恢复后补传记录；“多工程联合调度是否采用同一来水过程”属于流域协同问题，重点则是统一时空基准、模型版本和跨域交换协议。卡片让选型讨论围绕可验证行为展开，而不是围绕“平台大不大”争论。

当一个平台同时服务三种形态时，租户隔离、数据域和角色权限要在架构初期确定。流域分析员可以读取多个区域的汇总序列，工程值班员只能访问所属工程的原始状态，审批人可以查看影响范围和回退条件但不一定拥有设备操作权限。跨域调用要带调用方、用途、数据版本和有效期，服务端再次校验而不是信任前端传来的角色名。这样才能在共享代码和基础设施的同时，保持每个形态的责任边界清晰。

验收时还要按平台形态分别抽取样本。流域级抽查跨区域汇总和模型回放，区域级抽查计划、配水和统计口径，工程级抽查测点质量、值班确认和设备回执；三类样本都要能沿对象编码、时间基准和审计事件反向定位到原始记录。若同一水位在三个尺度显示不同，先检查聚合窗口、质量过滤和版本，而不是简单修改页面数字。抽查结果记录样本范围、观察时间和责任人，便于下一轮版本比较。

每次形态调整都应同步更新接口目录、权限矩阵、运行手册和回滚预案。这些交付物共同构成平台形态选择的可审计依据，也为后续容量、性能和安全评估提供共同输入；评估结论随版本和责任边界一并归档，归档包至少包含形态选择依据、数据域说明、接口版本、权限矩阵、样本结果和回滚负责人——这些材料同时是后续扩展到更大空间尺度时判断兼容性的基线。

\subsection{架构原则的“正例—反例—证据”三栏评审}

架构评审还可以采用“正例—反例—证据”三栏法。正例描述系统应当如何工作，反例描述一个看似方便但会造成责任或数据风险的实现，证据栏则指定测试数据、日志字段、权限记录或演练步骤。以“高风险操作受控”为例，正例是专业分析员提交调度建议、审批人批准、值班员按设备回执执行；反例是任何能打开页面的人都能调用控制接口；证据应包括角色令牌、审批事件、指令有效期、设备回执和回滚记录。以“全过程可观测”为例，正例是一次预警可以追溯到输入快照和模型版本，反例是只保存最终颜色，证据则是追踪ID、规则版本和复核结论。三栏法让原则变成可执行的评审任务。

每次架构决策都应形成简短记录：先写业务问题和受影响对象，再列出候选方案及其取舍，随后说明选择依据、风险、观察指标、回滚条件和责任人。记录不能只写“采用高可用架构”或“加强安全”，而要写成可检查的句子，例如“网络中断超过心跳窗口后，边缘节点继续保存原始观测并标记数据龄期；模型服务拒绝使用超过时效的输入；恢复后按事件ID去重并补写审计记录”。这样的决策记录既能指导实现，也能在维护阶段解释为什么某个看似保守的限制是必要的。

原则落地还需要跨角色复核。业务代表检查对象、岗位和处置流程，专业分析员检查单位、时间、质量码和模型边界，开发人员检查接口、异常和依赖，测试人员检查可复现判据，运维人员检查监控、备份和回退，安全人员检查身份、最小权限和审计留存。小团队可以由同一人兼任多个角色，但评审记录仍应分栏写明每个视角的结论。若某一条原则暂时不能验证，应标记为待补证据并给出关闭条件，不应以“后续再看”结束评审。

评审结论应绑定版本号和复核日期，避免同一条原则在后续迭代中失去上下文；每次复核保留参与角色和未关闭风险的清单，随发布包一并保存。


\section{生命周期选型评审记录与需求基线治理}\label{app:ext-ch02-baseline-governance}\label{online:ch02}

\paragraph{对应正文}第2章2.2、2.6节。

\subsection{选型评审记录与两周试验}

选型评审的输出至少包括：项目特征证据、模型组合、每个模型的前置条件、第一轮可验证交付物、主要风险、审批人和模型切换触发条件。项目进行中如果用户参与频率下降、接口风险暴露或合规范围扩大，应重新评估活动顺序、参与方式与验证条件，并更新模型选择记录。评审时可对照前后两版记录，说明调整依据及其对交付计划的影响。

实际评审可以采用“先排除、再组合、后试验”的顺序。先排除无法满足合规、现场条件或团队能力的模型；再把全局基线、局部原型、增量交付和风险轮次组合起来；最后用一个两周左右可观察的试验验证假设。试验不以代码行数为完成标准，而以是否获得关键证据为标准，例如接口是否能在断网后安全恢复、审批人能否独立完成退回与再提交、数据质量问题能否被追溯。试验结束后将结果写入选型记录，明确继续、调整或停止的决定，下一次评审直接引用这些证据。

评审记录还应标出参与者、评审日期和所依据的基线版本，下一轮评审先核对这些字段，确认各方比较的是同一范围，再讨论模型是否需要调整。记录中还要写明模型选择后的第一个可验证交付物、风险负责人、预期证据、观察窗口和停止条件；若证据不足，就延长验证或退回组合设计，不能用“进度已完成”替代模型有效性。人员、设备或试点范围变化后，可据此比较前后条件，核查原先的风险判断和验证结果是否仍然适用。

\subsection{需求基线的版本包、变更委员会与基线复核}

矩阵可以放在版本库中由构建检查，或由项目管理工具维护，但字段含义必须固定。需求ID、版本、来源、责任人、设计对象、接口路径、测试用例、验收结果和变更号共同关联需求来源、实现与验收结果；删除一条需求时保留废止原因和替代需求，不能让历史测试失去解释。发布前由需求负责人抽样反向追踪，确保“已完成”至少有一项可复现证据，未完成项明确进入下一版本。

实施变更时要维护“旧版本可读、新版本可验证、失败可回退”三条边界。数据库字段增加可以先向后兼容，再在新版本使用；字段含义改变时应新建版本或迁移脚本，保留旧字段读取期；接口响应新增字段要确认旧客户端忽略它不会破坏流程。涉及代理键和个人信息时，迁移脚本只搬运业务标识，不把身份证号等敏感字段复制到新主键；回滚时要检查审计记录和外部消息是否已经发出，不能只恢复数据库快照而留下无法解释的任务状态。

基线发布需要一个可复现的版本包。除了SRS正文，还应冻结引用的模型文件、数据字典、接口示例、测试数据、评审纪要和待确认问题；版本包生成校验摘要，发布人和审批人签名。开发、测试和运维使用同一摘要确认自己拿到的是同一份需求，遇到争议时先核对版本再讨论实现。修订时保留差异说明，指出新增、修改、废止和迁移字段，不能用“全文更新”四个字替代影响分析。

变更委员会的职责可以按四类角色分工：业务代表确认价值和范围，安全与合规人员确认硬约束，技术负责人确认实现和迁移风险，测试与运维负责人确认验证、观察和回退条件。小团队可以由同一人兼任多个角色，但每个角色的决策记录仍要单独填写。涉及闸门控制或防洪调度的变更，即使只是调整默认参数，也应至少经过审批人和专业分析员复核；普通报表排序可以采用轻量流程，但仍要留下版本号和回归用例。

一份合格的SRS还应在交付时附上阅读说明：先读范围和角色，再读功能用例与数据契约，最后对照质量指标、追踪矩阵和变更历史。读者按照同一路径检查，能在不依赖作者口头解释的情况下复现关键决策和验收步骤。

需求基线发布后还要安排一次“基线复核”，确认文档、模型、接口示例和测试数据来自同一版本摘要。复核时随机抽取几条需求，从业务目标追到验收证据，再从测试失败反查责任人和变更号；如果任一方向出现断点，就把断点登记为发布风险并规定关闭条件。对案例灌区这类持续接收设备数据的系统，复核还应覆盖跨日时段、重复提交、权限变化和网络恢复等边界场景，确保文字规则在真实运行节奏下仍然可执行。只有通过基线复核，版本包才适合作为开发、测试、培训和运维共同引用的唯一依据。

\section{第3章补充阅读：需求复述与原型方法}\label{online:ch03}

\paragraph{对应正文}第3章3.1.1节（需求与用例）与3.2.3节（用原型检验分层和模块的划分）。正文保留 REQ-MON 需求登记和三个核心用例、架构原型的定义；本节收录多终端、可靠性、安全与易用性需求的复述，以及原型构建、反馈收集与细化的通用方法。

\subsection{平台的非功能需求复述}

平台还应具备多终端协同能力。值班员希望通过监控大屏实时掌握水库运行态势，通过 Web 端查看分测点趋势曲线和告警记录，通过移动端在巡检现场快速核对设备状态和处理结果。专业分析员需要对预警阈值、分析模型和应急流程进行研判，运维员则要完成设备注册、通信链路维护、账号权限控制以及与上级监管平台的数据交换。

系统的可靠性和实时性是水利业务场景中的首要非功能需求。监测平台需要在汛期、强降雨、突发险情等高压工况下保持持续运行，即使出现局部断电、边缘站点离线或通信网络抖动，也要通过断点续传、缓存补报、双链路通信和主备切换等机制保证重要数据不丢失、预警链路不中断。例如，前置采集网关可在通信恢复后自动补传数据，平台中心可对关键告警采用短信、应用消息和声光告警等多通道并发推送。

平台承载工程安全数据和调度决策信息，泄露、篡改或越权操作可能直接影响研判与处置，因此需要在采集链路、平台服务、用户访问和审计追踪等层面建立防护机制。数据传输过程中应采用 TLS、VPN 或专网加密，重要操作应保留审计日志，重要告警应支持签名和确认回执，用户认证应结合统一身份认证和分级授权策略，确保监测数据、预警结论和处置指令不会被篡改或误用。

平台的易用性与可扩展性也需要一并考虑。监测界面应突出重点测点、重要告警和风险趋势，避免在紧急场景下出现信息淹没；系统架构则应支持后续接入新的监测设备、新增分析模型或扩展到河道、泵站、灌区等更多水利场景。

\subsection{原型的构建与细化}

\paragraph{原型的定义与作用}

原型是软件系统的初步版本，它具有真实系统的部分功能和特性，但通常在完整性和性能上不如最终产品。对于水利工程安全监测平台而言，原型可以是一个简化版本，具备基本的监测数据接入、简单的预警功能以及基础的用户交互界面。其主要作用在于：

\paragraph{直观呈现设计理念}通过原型，项目团队能够将抽象的架构设计转化为具体可操作的实体，让利益相关者（如业务部门、开发团队成员、管理层等）直观地感受到系统的功能和使用方式。这有助于各方更好地理解架构设计，发现潜在的需求不一致或设计缺陷。例如，值班员可以通过操作原型，提前发现某些功能不符合实际使用习惯，及时提出修改意见，避免在开发后期进行大规模返工。

\paragraph{验证架构的可行性}在水利工程安全监测平台的开发中，通过构建原型，可以验证系统架构在技术上是否可行。比如，测试在模拟多个监测点同时工作的情况下，系统的数据处理和传输是否能够满足实时性要求；验证系统各模块之间的通信和协作是否顺畅，是否存在接口不兼容或交互异常的问题。如果在原型阶段发现问题，可以及时调整架构设计，降低后期开发风险和成本。

\paragraph{促进沟通与协作}原型为项目团队成员、业务部门以及其他利益相关者提供了一个共同的沟通平台。在讨论和评估原型的过程中，各方可以充分交流意见和想法，明确项目的目标和需求。开发团队可以根据业务反馈，进一步细化和优化架构设计，提高项目成功率。

\paragraph{快速原型的构建方法}

快速原型的构建需要采用一些高效的方法和工具，以快速实现原型并进行验证。

\paragraph{选择合适的技术栈}根据项目需求和团队能力，选择能够快速开发的技术框架和工具。例如，对于水利工程安全监测平台的前端原型，可以使用 HTML、CSS 和 JavaScript，结合 Vue.js 等框架快速搭建监测总览界面；后端可以选择轻量级框架快速实现数据处理和接口功能。利用模拟数据生成工具为原型提供监测数据和业务数据，以便进行功能测试。

\paragraph{采用迭代式开发}快速原型通常采用迭代式的开发方式。在每个迭代周期中，逐步增加原型的功能和完善其细节。例如，在第一个迭代中，实现基本的传感器数据显示功能；在第二个迭代中，添加预警功能的初步实现；在后续迭代中，不断优化用户界面和系统性能。通过迭代式开发，可以及时根据反馈调整原型，使其逐渐接近最终产品的要求。

\paragraph{注重核心功能的实现}在构建原型时，应重点关注系统的核心功能和业务流程。对于水利工程安全监测平台来说，监测数据接入、异常预警和用户查询等是核心功能。优先实现这些功能可以快速验证架构的可行性。对于一些辅助功能和细节，可以在后续优化阶段完善，避免在原型阶段花费过多时间和精力。

\paragraph{原型的细化与优化}

在完成初步的原型构建后，需要根据反馈对原型进行细化和优化。

\paragraph{收集反馈意见}通过组织利益相关者对原型进行试用和评估，收集各方的反馈意见。这些意见可能包括功能需求的调整、用户界面的改进建议、性能优化的要求等。例如，用户可能认为原型的操作流程不够便捷，需要简化；开发团队可能发现某些功能在技术实现上存在困难，需要调整架构设计。

在收集反馈时，应采用多样化的方法，如用户测试、问卷调查、小组讨论等，以确保获取全面且有价值的信息。对于水利工程安全监测平台，可以组织值班员和专业分析员使用原型，观察他们在监测查询、告警确认和图表分析过程中的操作行为和遇到的问题。向利益相关者发放详细问卷，涵盖功能完整性、易用性和性能表现等多个方面，并组织开发团队、产品负责人和业务代表进行讨论，共同分析问题产生的原因和可能的解决方案。

\paragraph{功能细化与扩展}根据反馈意见，对原型的功能进行细化和扩展。比如，在水利工程安全监测平台原型中，最初的预警功能可能只是简单弹窗提示，在细化阶段，可以增加告警推送功能，将预警信息及时发送到值班员的终端，并提供监测点位置、指标类型和趋势变化等信息。

进一步扩展水利工程安全监测平台的预警功能，除了向值班员终端推送预警信息外，还可以与调度平台或应急值守系统进行联动。当检测到渗流异常或位移超限时，通知值班员，也能自动将相关信息发送到调度中心，包括工程位置、异常指标、发生时间等详细信息，以便快速响应。为预警功能增加智能分类能力，根据监测数据和预设规则，区分渗流异常、变形异常、水位异常等不同类型，针对不同类型的预警提供差异化处置提示。

\paragraph{性能优化}对原型的性能进行优化，确保在实际使用场景下能够满足性能要求。例如，优化传感器数据的处理算法，提高数据处理的速度和准确性；优化系统的网络通信机制，减少数据传输的延迟。

性能优化应从测量结果出发。处理耗时较高时，先比较过滤与分析算法；若试用机器学习识别异常趋势，应在同一组标注样本上测量漏报、误报和推理耗时。网络受限时，可测试压缩与异步传输，并检查压缩开销和队列等待是否抵消收益；存储受限时，先检查查询与索引，再评估时序存储的扩展方案；并发请求增加时，可评估负载均衡。每次调整使用相同输入与负载，记录时延和资源占用，发现性能退化时恢复上一配置。

\section{服务边界、接口版本与健康检查}\label{app:ext-service-boundary}\label{online:ch05-api}

\paragraph{对应正文}第5章5.1节与第8章8.6节。

案例水库采用模块化单体（第3章\ref{sec:ch03-load-estimate}节）。下面的内容在平台拆成多个服务、或者接口对外公开之后才会用到。

服务边界按业务能力划分。数据采集、观测查询、预警评估和通知可以是不同的服务，每个服务有明确的数据所有权和 API 契约；拆分不是把每个类都部署成独立进程。开发环境用配置文件把逻辑服务名映射到固定地址，生产环境再接服务注册与发现、健康检查和负载均衡。业务代码只依赖接口客户端，注册中心的注解不进领域层，基础设施和业务规则才能分别修改和测试。

跨服务调用显式处理超时、重试和降级。查测站详情可以在短超时后返回缓存快照，界面标注数据时间；写观测值不能因为重试而重复落库，请求 ID 和幂等键从网关一直带到服务日志和数据库。链路追踪传\texttt{traceId}和\texttt{requestId}，错误应答给可定位的错误码，不暴露内部堆栈。两个服务共享同一张 PostgreSQL 表，边界就名存实亡了，改成通过版本化 API 或消息事件交换数据。

接口演进要兼容，也要留记录。对外公开的接口用路径版本\texttt{/api/v1/assets}，加字段时旧字段语义不变；删字段、改单位、改时区都是破坏性变更，发布\texttt{v2}并给迁移窗口。版本号不随每次修复递增，补丁和兼容的字段增加写进文档和变更日志。应答里带\texttt{Deprecation}和\texttt{Sunset}头，调用方在旧版本停用前有时间升级。

服务发现的健康检查分进程存活、依赖可达和业务可用三类，含义与第8章8.6.1节相同：进程存活用来重启异常实例，依赖可达检查数据库、Redis 和 Kafka 连接，业务可用验证关键表结构、时序写入权限和预警规则加载。网关切换路由时先摘掉不健康的实例，再等正在处理的请求完成；客户端只对明确的瞬态错误重试，写入请求带幂等键。


游标分页适合持续追加的时序读数：客户端携带上一页最后一条记录的\texttt{occurredAt,id}，服务端按联合索引向后查找，避免高页码\texttt{OFFSET}扫描大量历史数据。游标应视为不透明字符串，不能让客户端自行修改时间边界；接口文档还要说明读数在分页期间追加时的排序稳定性。报表导出属于异步任务时，使用任务资源返回\texttt{202 Accepted}，由客户端轮询任务状态，而不是让一次 HTTP 请求长时间占用连接。

接口文档还应给出示例请求、示例响应、错误码和限流规则，并在每次版本发布时同步更新 OpenAPI 描述；值班员和前端开发者据此可以复现问题、核对字段单位并确认兼容范围。

\section{坐标成果的入库、审计与版本管理}\label{app:ext-ch06-crs-audit}\label{online:ch06-crs}

\paragraph{对应正文}第6章6.2节。

第6章6.2.1至6.2.3节给出了坐标、高程和局部原点的换算。数据来自多个单位、平台长期运行时，还要把换算过程管起来。

\paragraph{元数据分两级保存}数据集级元数据保存统一的坐标系和版本：水平坐标系、投影方法、带宽、带号、中央经线、尺度因子、假东与假北、东坐标是否带带号前缀、高程类型、高程基准、单位、精度等级、转换模型版本和生产软件版本。要素级属性保存测站、控制点或构件的采集方法、精度和质量标识。同一图层混有不同来源的数据时，入库时拆成多个数据集并分别建立转换记录；来自外单位的数据同时保存原始坐标、转换前后的 EPSG 代码和转换日期。转换成功不等于精度合格，转换记录里还要有控制点数量、残差统计和审核人，成果复核时可以重算。

\paragraph{原始区、标准区、发布区}测量成果进入平台后分三层存放。原始区保存外业文件和原始坐标，不覆盖；标准区完成 EPSG、单位、高程基准和质量码的统一，生成带版本的数据库表；发布区按 WMS、WFS 或三维瓦片的需要生成切片、索引和缓存。每层保存源文件哈希、处理工具、参数和责任人，服务刷新只读取通过质量检查的标准区数据。出现定位问题时，可以按哈希和版本回放整条处理链。

\paragraph{转换流水线的检查点}接收阶段读取坐标系和高程基准，缺少元数据的数据进入隔离区，补全后重新校验；预处理阶段按带宽、中央经线和轴序完成转换，保留原始文件和工具版本；质量阶段用独立控制点计算残差，平面和高程分开统计最大绝对误差、均方根误差和样本数，检查是否存在整体平移、旋转或比例误差；发布阶段生成服务能力文档、EPSG 声明和数据字典；运行阶段新增测点继承同一套转换配置，配置变更时重新运行回归检查。控制点要覆盖坝轴线两端、闸室、廊道出入口和场景边界，不能集中在一小块区域。

\paragraph{精度预算与变更控制}坐标精度预算在需求阶段分配到数据采集、控制测量、转换计算、模型简化和渲染显示各环节。测量层的毫米级精度经过投影和局部原点平移后仍要可逆；模型简化允许视觉误差，但不改变用于定位闸门或测点的控制顶点；屏幕显示可以按像素取整，查询和分析仍使用原始的米制坐标。每次改变中央经线、局部原点或高程模型，重新生成控制点报告，并与上一版本比较残差和服务范围。

长期运行的平台把坐标配置当作基础设施契约，而不是前端常量。配置中心记录当前版本，后端接口在应答里返回坐标版本，前端加载模型时核对版本；不匹配时显示待更新状态并停止自动叠加，避免把新测量数据投到旧场景上。坐标版本更新要通知地图服务、三维瓦片和监测接口的维护者，发布说明列出受影响的图层、测站和查询时间窗，值班员据此抽查。

\paragraph{可重复性测试}用同一组经纬度和元数据在开发机、持续集成环境和生产镜像里各转换一次，比较东、北坐标和高程；差异超过约定精度时冻结发布，检查 PROJ 数据文件、浮点模式和参数顺序。测试样本既包含中央经线附近的点，也包含分带边缘的点和高程异常符号变化的点。部署后抽查一个测站，依次执行地图点选、三维拾取、后端查询和原始坐标回放，核对各步骤指向同一对象。验收记录写明转换日期、操作者、软件环境以及输入、输出文件的哈希；更新 PROJ 数据库或高程模型时，先在隔离环境重算并比较差异，旧版本数据保持可查询。

\paragraph{教学实验}先用一个已知控制点验证公式，再扩展到整批测站。实验报告列出原始经度、带宽、带号、中央经线、投影参数、转换后坐标、局部原点和误差统计；使用软件默认值时说明默认值的来源和适用条件。再做一个故意改变轴序或高程类型的反例，观察地图、三维模型和查询结果怎样分离。


\section{逐句追加的复核要求汇编}\label{online:review-notes}

\paragraph{对应正文}第6章6.2.4节（OGC地图与要素服务）、第7章7.4.3节（触控、长按与世界单位）与第8章8.3.6节（前端：预警状态与处置操作）。这些句子原以单句逐行追加在段落末尾，内容为复核与归档的细节要求，现集中收录于此。

对异常响应还要检查Content-Type、错误码和重试建议，前端按契约展示可操作的提示。服务端同时限制错误详情长度，避免把数据库或文件系统信息暴露给外部用户。客户端日志只保留必要上下文和traceId。故障恢复后再开放自动刷新。恢复记录应关联演练编号和验证人员。并标注验证结果。记录已归档并可追溯。后续复测沿用同一编号和样本。

回归报告还应标出渲染器、模型块版本和浏览器版本，使同一坐标偏差能够在软件升级后定位和比较。同时记录设备像素比和当前LOD级别，便于比较不同屏幕下的命中范围与帧时间。还要保留测试输入和预期对象清单，确保回归结果可由其他同学独立复现。提交前再次复核截图与日志。

页面还要在错误恢复后保留原请求的追踪号和用户选择，便于值班员确认恢复是否真的生效。验收记录同时保留浏览器版本、构建版本和接口版本，便于复现跨端差异。

\section{第8章补充阅读：参数驱动的测试与审计}\label{online:ch08-params}

\paragraph{对应正文}第8章8.4.3节（模型可信度与人在回路）。正文保留特征水位边界点的测试覆盖与版本化参数集两段；本节收录追踪号贯通、前端校验提示与降雨情景解释等内容。

模型运行、闸门反算和水网平衡共享同一追踪号。追踪号串起输入水位、闸门开度、公式版本、输出流量、未满足需求、闭合残差和审批回执。这样即使后来替换了水力模型，审计人员仍可依据旧版本重现当时的计算，而不是用最新模型解释历史处置。输入快照与公式版本共同确定一次计算的复现条件。
工程参数还要进入接口和前端的校验提示。用户输入闸门开度时，页面显示来自版本化参数集的最小值、最大值、单位和生效时间；输入水位时，页面同时显示死水位、汛限水位、正常蓄水位、设计洪水位和校核洪水位，避免只给一个“超限”布尔值。接口拒绝越界参数时返回字段路径、参数版本和建议动作，调用方可以把错误回显给专业人员，而不是把后端异常堆栈直接展示给值班员。测试夹具固定三类边界：刚好等于阈值、略低于阈值和略高于阈值，并验证蓝黄橙红分级、闸门公式适用性和水量闭合残差在同一版本下保持一致。

对于降雨情景，平台把流域面积320\,km$^2$作为输入范围与结果解释的背景，不把面积直接当作瞬时来水量。三种降雨情景共享同一初始水位和闸门设备参数，只改变降雨过程与不确定性说明；输出的水位包络线标注情景名、时间步长和模型版本。若情景结果超过校核洪水位，系统先生成会商任务，再允许专业人员比较提前泄洪、维持水位和分时调度等方案。情景、参数和批准方案按版本关联。每次发布还应把参数快照和计算结果校验和写入审计日志，供恢复演练与课程验收复核。

\section{基于构件的设计与复用}\label{app:ext-component-based-design}

\paragraph{对应正文}第3章相关小节，具体入口见下文说明。

第3章3.4节按“因为什么而改动”划分模块，并用一个具体变化检验边界。模块如果要跨项目复用，例如把案例水库的监测接入、认证与审计用到灌区或河道项目上，还需要多做三件事：从领域里找出通用构件，接入前检验构件是否合格，把接口和配置设计成适合复用的样子。这套做法称为基于构件的设计（Component-Based Design，CBD），供课程设计选读。

\subsection{领域工程：找出可复用的构件}

领域工程是对某一类应用做调研，找出其中反复出现的通用功能，把它们做成可复用的构件。在水利信息系统里，身份验证、日志与审计、设备接入、观测数据的质量检查几乎每个项目都需要，流程和规则也大同小异，适合提取。

提取分两步。第一步识别通用功能：比较领域内的几个应用，列出它们共有的功能模块。第二步设计通用构件并封装领域知识：例如通用的身份验证构件除了校验用户名和密码，还把用户管理、权限控制的相关规则一并封装在内部；新项目只需配置和调用，不用重新开发。

\subsection{构件合格性检验}

构件接入系统之前，要检验它的接口行为、性能和质量约束是否满足需求。检验结果决定三种去向：直接采用，增加一层适配后采用，或者更换实现。

\paragraph{单元测试}单元测试验证构件内部的逻辑。对一个数据处理构件，检查它对各种输入的转换和计算是否正确，及时发现并修复内部缺陷。

\paragraph{集成测试}集成测试在构件装进系统以后进行，检查它与其他模块的交互是否正常。以通信构件为例，重点观察三种失败形态：连接建立不上；报文能发出但字段被截断；重传后同一条读数被写入两次。这三种情况在单元测试里都看不出来，只有把构件放回真实的系统环境才会暴露。

\subsection{面向复用的接口与配置设计}

\paragraph{接口标准化}对外提供通用的接口形式（如 REST API），调用方就不必为每个新系统改写一次适配代码。监测数据接入构件如果提供标准化的 REST 接口，监测平台、移动巡检应用和专题分析系统都可以直接调用，不用关心它的内部实现。

\paragraph{配置灵活性}把随场景变化的部分做成参数。传感器数据采集构件可以通过配置文件指定传感器类型、连接方式和采样频率，同一个构件就能适应不同类型的传感器和不同的工作环境。

\paragraph{案例}案例水库的监测接入模块遵循通用的接口标准和通信协议，不依赖特定业务系统的内部实现，因此可以移植到其他水利信息系统：河道水文监测系统用它采集流量、水位和雨量，泵站运行监测系统用它接入振动、电流和视频数据，各自更换一批协议适配器即可。



\section{对话框焦点管理与长列表渲染}\label{app:ext-fe-focus-list}

\paragraph{对应正文}第4章相关小节，具体入口见下文说明。

本节两段脚本接第4章4.2.3节和4.4.6节，需要4.4节的 DOM 与事件基础。

\paragraph{对话框的焦点管理}弹出详情对话框时，键盘焦点要移到对话框标题；关闭时回到触发它的按钮，值班员才能接着刚才的位置继续操作。清单\ref{lst:ch04-a41-focus}用原生\texttt{<dialog>}元素实现：\texttt{openDialog}先记下当前焦点所在的元素，\texttt{showModal()}打开模态对话框后把焦点交给带\texttt{autofocus}的标题；\texttt{closeDialog}关闭对话框并把焦点还回去；按 Esc 触发的\texttt{cancel}事件走同一个关闭函数。脚本只管交互，样式仍交给 CSS。换成 Vue 组件时，把两个函数放进组件，并在\texttt{onUnmounted}里移除事件监听。在控制台运行后只用键盘操作：回车打开，Esc 关闭，焦点应回到“查看测站详情”按钮上。

\begin{lstlisting}[label={lst:ch04-a41-focus},language=JavaScript,caption={监测详情对话框的焦点管理}]
document.body.insertAdjacentHTML('beforeend', `
  <button id="open-station" type="button">查看测站详情</button>
  <dialog id="station-dialog">
    <h2 autofocus>测站详情</h2>
    <p>此处显示经过后端校验的测站信息。</p>
    <button data-close type="button">关闭</button>
  </dialog>`);

const dialog = document.querySelector('#station-dialog');
const openButton = document.querySelector('#open-station');
let lastFocused = null;

function openDialog() {
  lastFocused = document.activeElement;
  dialog.showModal();
  dialog.querySelector('[autofocus]')?.focus();
}

function closeDialog() {
  dialog.close();
  lastFocused?.focus();
}

openButton.addEventListener('click', openDialog);
dialog.querySelector('[data-close]')
  .addEventListener('click', closeDialog);
dialog.addEventListener('cancel', closeDialog);
\end{lstlisting}

\paragraph{只渲染可视窗口的长列表}列表到上万行时，瓶颈在节点数量和布局计算。虚拟滚动只为可视区域附近的几行创建节点，并把滚动位置换算成数据下标。清单\ref{lst:ch04-a44-virtual-list}是固定行高时的基本思路：占位层把滚动条撑到全量高度，可见的 8 行用绝对定位放到各自的位置，滚动时重新计算起始下标。运行后拖动滚动条，“元素”面板里始终只有 8 个行节点。实际使用还要处理行高不固定、键盘焦点落在被回收的行上，以及向读屏软件说明“第几项，共几项”；测点卡片高度差别很大时，先用分页或分组折叠，确认性能不够再引入虚拟滚动。

\begin{lstlisting}[label={lst:ch04-a44-virtual-list},language=JavaScript,caption={长测站列表的可视窗口渲染思路}]
const root = document.createElement('div');
root.innerHTML = '<div id="virtual-list" style="height:160px;overflow:auto"></div>';
document.body.append(root);
const viewport = root.querySelector('#virtual-list');
const records = Array.from({ length: 1000 }, (_, index) => ({
  assetId: `DAM-A-PZ-${String(index + 1).padStart(4, '0')}`,
  value: (180 + index / 100).toFixed(2)
}));
const rowHeight = 32;
// 占位层把滚动条撑到全量高度，可见行再absolute定位到对应位置；
// 没有它 scrollTop 永远到不了后面的行
const spacer = document.createElement('div');
spacer.style.cssText =
  `position:relative;height:${records.length * rowHeight}px`;
viewport.append(spacer);

function renderWindow() {
  const first = Math.floor(viewport.scrollTop / rowHeight);
  const visible = records.slice(first, first + 8);
  spacer.replaceChildren(...visible.map((record, offset) => {
    const row = document.createElement('div');
    row.style.cssText = `position:absolute;left:0;right:0;` +
      `top:${(first + offset) * rowHeight}px;height:${rowHeight}px`;
    row.textContent = `${record.assetId}：${record.value} kPa`;
    return row;
  }));
}
viewport.addEventListener('scroll', renderWindow, { passive: true });
renderWindow();
\end{lstlisting}



\section{明暗主题与高对比度样式}\label{app:ext-fe-theme}

\paragraph{对应正文}第4章相关小节，具体入口见下文说明。

CSS 自定义属性在运行时可以被继承和覆盖，适合承载一整套主题。文字色、背景色、边框色、状态色和间距集中定义在\texttt{:root}或页面容器上，组件通过\texttt{var()}读取；值班室降低环境亮度时，只替换这一组变量，地图、卡片和表格保持同样的对比关系。暗色主题不能简单地把白色换成黑色，还要重新检查文字与背景的对比度、曲线颜色的区分度、预警颜色在暗底上的可见性和焦点环的边界。

主题由用户通过按钮或系统偏好决定，不随业务预警等级自动切换。\texttt{prefers-color-scheme}作为初始值，用户显式选择后存到本地并覆盖系统偏好；主题切换只改变呈现用的变量，不改变 DOM 顺序和语义。截图、打印和投影场景还要有高对比度的后备色，“正常”与“预警”不能只靠颜色区分。

清单\ref{lst:ch04-a43-dark-theme}给出明暗主题和高对比度状态样式。\texttt{data-theme}属性由应用状态控制；\texttt{forced-colors}分支让操作系统的高对比度模式接管边框和文本颜色。把这段样式与清单\ref{lst:ch04-a41-landmarks}的页面组合后，先用窄屏和键盘完成基本任务，再到 1920、2560 和 3840 像素的画布上检查，最后切换暗色和高对比度模式，每种条件下水位、预警和时间信息都应可读。

\begin{lstlisting}[label={lst:ch04-a43-dark-theme},language=CSS,caption={监测页面的明暗主题与高对比度}]
:root {
  --surface: #ffffff;
  --surface-raised: #f5f9ff;
  --ink: #263238;
  --border: #b0bec5;
  --normal: #1565c0;
  --warning: #d46b08;
}

[data-theme="dark"] {
  --surface: #17212b;
  --surface-raised: #0f151b;
  --ink: #f5f7fa;
  --border: #607d8b;
  --normal: #90caf9;
  --warning: #ffb74d;
}

@media (prefers-color-scheme: dark) {
  :root:not([data-theme="light"]) { color-scheme: dark; }
}

.monitor-panel {
  color: var(--ink);
  background: var(--surface);
  border: 1px solid var(--border);
}
.monitor-panel.is-warning {
  outline: .2rem solid var(--warning);
}
@media (forced-colors: active) {
  .monitor-panel.is-warning { outline: 2px solid CanvasText; }
}
\end{lstlisting}



\section{前端构建产物的发布与回滚}\label{app:ext-fe-release}

\paragraph{对应正文}第4章相关小节，具体入口见下文说明。

本节承接第4章4.8.3节。清单\ref{lst:ch04-a48-deploy}给出构建与部署的命令和反向代理的两条关键规则：前端容器只提供静态文件，\texttt{/api/}转给后端，其余找不到的路径回退到\texttt{index.html}。

\begin{lstlisting}[label={lst:ch04-a48-deploy},language=bash,caption={Vite 产物与 Docker Compose 部署契约}]
# 构建阶段：在固定 Node 版本中生成 dist/
docker build --target build -t water-web-build .
# 运行阶段：静态服务器只读挂载产物，API 由反向代理转发
docker compose up -d water-web nginx

# nginx 关键规则（伪配置，部署时写入站点配置）
location /api/ { proxy_pass http://water-api:8080/; }
location / { try_files $uri $uri/ /index.html; }
\end{lstlisting}

\paragraph{深链接与缓存分层}浏览器直接打开\texttt{/assets/DAM-A-PZ-07}时，静态服务器先找同名文件，找不到就回退到\texttt{index.html}，再由 Vue Router 渲染详情页；\texttt{/api/}必须先于这条回退规则匹配。缓存头分三层：\texttt{index.html}短缓存并要求重新验证，避免用户拿到旧入口去请求已经删除的代码块；带内容哈希的脚本和样式长期缓存；监测接口按数据时效禁止缓存或设置很短的有效期。部署后用真实浏览器和命令行各验证一次，只看服务器状态码不够。

\paragraph{构建流水线}流水线分为依赖安装、静态检查、单元测试、生产构建、产物检查和发布六步。依赖安装使用锁文件（\texttt{npm ci}），不同机器得到相同版本；静态检查关注未使用变量、可访问性和路径别名解析；测试至少覆盖路由守卫、筛选计算和请求错误分类；构建日志记录 Node、Vite 和插件版本；产物体积超过预算时阻断发布。

\paragraph{不可变目录与回滚}每次构建把\texttt{dist/}复制为带提交号的版本目录，生成一份包含入口哈希、资源清单和构建时间的 manifest；反向代理通过符号链接或配置项指向当前版本。切换后依次检查首页、登录、测点详情和接口健康状态。发现路由 404、资源加载失败或接口异常时，切回上一个 manifest 即可恢复。回滚只改变静态文件的指向，不动数据库和时序数据；上一个版本保留到确认监测业务稳定之后再清理。

\paragraph{产物安全检查}检查分源代码、产物和运行时三个面。源代码扫描禁止把令牌、私钥和内部密码提交到仓库；产物扫描检查压缩后的 JavaScript 是否残留调试开关、源代码地图是否暴露内部路径；运行时检查内容安全策略、HTTPS、跨域策略和错误页面是否泄露堆栈。环境变量的\texttt{VITE\_}前缀不是保密机制，进入前端产物的值一律视为公开。发布记录写明发布人、构建提交号、镜像摘要、变更范围、接口兼容性、验证结果和回滚入口，值班员交接班时据此判断当前版本状态。



\section{地图服务的发布、缓存与运维}\label{app:ext-ch06-gis-ops}

\paragraph{对应正文}第6章相关小节，具体入口见下文说明。

第6章6.2.4节说明了 WMS、WMTS、WFS、WCS 和 3D Tiles 各返回什么，以及怎样在 Cesium 里接入自建底图。把这些服务长期运行起来，还要处理发布、缓存、版本和故障。

\paragraph{各类服务的数据生命周期}WMS 的样式和图例变化时重新渲染即可，不必重建原始数据；WMTS 适合版本稳定的底图，更新时按区域和层级增量切片；WFS 的属性字段属于业务契约，字段改名通过版本化接口发布；WCS 的覆盖数据保留像元分辨率、NoData 值、单位和采样时间，供模型计算复现；3D Tiles 在每个瓦片里维护几何、纹理和层级元数据，客户端按视域和屏幕误差请求。这几类服务的缓存策略、权限边界和更新节奏各不相同，不宜包装成一个笼统的“地图接口”。

\paragraph{GeoServer 发布步骤}发布分四步，每一步都可以单独回退：在工作区登记数据源和坐标元数据；创建工作空间和图层；配置样式与访问权限；在 GeoWebCache（GWC）里创建或绑定瓦片矩阵集（gridset）并预生成热点区域的缓存。发布前用 GetCapabilities 确认服务版本、图层名和坐标系，发布后用一条 GetMap、一条 GetFeature 和一条 GetCoverage 请求做冒烟测试：GetMap 看颜色和图例，GetFeature 查属性、几何和轴序，GetCoverage 抽几个像元与原始文件比对，3D Tiles 检查 LOD 切换、纹理色彩空间和业务标识是否保留。冒烟通过后再启动预生成；预生成任务按层级和工程范围拆分，失败的任务可以重跑，缓存目录带版本号，避免新旧瓦片混用。缓存清理和重新切片要有任务编号与进度记录，不在高峰时段整体删除。

GeoServer 的 GWC 默认只内置 EPSG:4326 与 EPSG:900913 两个瓦片矩阵集，CGCS2000 经纬度（EPSG:4490）的需要手工创建，并核对原点、瓦片宽高、分辨率和层级数。清单\ref{lst:ch06-geoserver-gridset}是一个教学配置片段：层0让$1^\circ\times1^\circ$的范围恰好落入一张$256\times256$的瓦片，以后逐层减半。

\begin{lstlisting}[label={lst:ch06-geoserver-gridset},language=YAML,breaklines=true,caption={GeoServer GWC EPSG:4490 gridset配置要点}]
gridSet:
  name: EPSG:4490
  srs: 4490
  tileWidth: 256
  tileHeight: 256
  extent: [113.0, 34.0, 114.0, 35.0]
  metersPerUnit: 111319.490793
  levels:
    # 层 0 让 1°×1° 范围恰好落入一张 256×256 瓦片：
    # 1/256 ≈ 0.00390625 度/像素，逐层减半
    - {level: 0, resolution: 0.00390625}
    - {level: 1, resolution: 0.001953125}
  cacheLayers:
    - water:cgcs2000_basemap
\end{lstlisting}

实际工程从测区范围和服务规范计算分辨率、矩阵宽高与最大层级。图层声明的坐标系与瓦片矩阵集不一致时，缓存在切片阶段就会错位；最大层级没有与 Cesium 的\texttt{maximumLevel}同步时，客户端会持续请求不存在的瓦片，日志里出现大量404。发布脚本把这两项作为启动前检查。

\paragraph{能力文档、缓存键与监控}每个服务维护一份能力文档，列出服务版本、操作、图层或覆盖名称、坐标系、轴序、输出格式、最大请求范围和异常编码，示例请求使用脱敏域名和小范围数据。样式变更会改变 WMS 的像素结果，但不改变 WFS 属性和 WCS 像元，所以图层发布记录分别保存样式版本和数据版本。缓存键至少包含图层、样式、坐标系、范围、宽高、时间维度和数据版本；带时间维度的水位或降雨图层，时间参数必须进入缓存键，否则不同时刻会显示同一张图片。平台启动时对 GetCapabilities 做健康检查，运行中记录应答时延、返回字节数、缓存命中率、4xx 与 5xx 数量和上游数据库耗时。WMS 渲染失败时，前端提示专题图不可用并保留基础底图；WFS 或 WCS 不可用时，相关查询按钮置为不可用，不把空结果显示成“没有测站”或“没有积水”。查询参数进入数据库或栅格处理器之前做白名单校验，日志记录请求摘要、追踪号、图层、坐标系和耗时，不记录完整令牌。

\paragraph{3D Tiles 发布检查}检查瓦片内容的空间参考、几何误差和属性透传。每个瓦片的包围体用于视锥裁剪和屏幕空间误差判断，包围体错误会造成过早卸载或加载过多；业务属性通过稳定的对象编码关联到后端，不把完整的监测记录塞进瓦片。转换工具的输出日志包含源模型、坐标转换、LOD 阈值、纹理压缩和失败对象，抽样加载根节点、中间层和叶节点之后再开放服务。模型更新使用新目录和新版本号，客户端完成切换后再回收旧目录，避免长连接用户读到发布了一半的状态。

\paragraph{向 OGC API 迁移}OGC API 路线强调资源、链接和可发现性：Features 返回集合与分页链接，Tiles 返回瓦片集和地址模板，Coverages 返回覆盖描述与取子集的能力，Maps 返回渲染结果及样式信息。经典 WMS、WFS、WCS 仍是存量系统的主流，平台可以采用“兼容层加新接口”的做法：内部数据模型统一，对外按客户端能力提供两种接口；代理层把旧的键值对参数转换为新的路径和查询参数，并记录新旧请求的对应关系。迁移期间对同一测站和同一块 DEM 样本做双接口比对，坐标、属性数量和像元统计一致后再逐步切换；能力文档标出推荐接口和弃用时间，监控按接口版本分别统计流量，确认没有客户端被意外切断。

\paragraph{故障演练}人为关闭 GWC、缩小 WCS 允许范围或撤销某个图层的权限，观察客户端是否显示清楚的状态、服务端是否留下审计记录、缓存是否保持版本隔离。再用错误坐标系、越界范围、超大图片、未知图层和过期令牌等异常输入逐一请求，检查服务是否返回稳定的错误码和正确的 Content-Type，错误详情是否泄露内部路径、数据库或文件系统信息。演练结果写入运行手册，下一次发布前按同一脚本、同一批样本复核。



\section{倾斜摄影的外业设计、空三检核与成果验收}\label{app:ext-ch06-photogrammetry}

\paragraph{对应正文}第6章相关小节，具体入口见下文说明。

第6章6.3.1节说明了倾斜摄影的流程、四种点和精度评定公式。下面是生产环节的细则，供参加实际项目的读者查阅；各项限差以项目适用的规范条款为准。

\paragraph{外业设计}影像覆盖范围要超过成果边界，航带端部和转弯区留出缓冲；建筑物立面、坝肩和峡谷等遮挡严重的部位增加交叉航线或补拍方向，并在任务书里记录补拍原因。曝光时间、光圈和感光度与飞行速度、光照和地表反射率匹配，避免运动模糊和高光饱和。飞行日志保存航线版本、起降时间、设备编号、镜头组合、定位定姿数据（POS）来源、坐标基准、天气和异常处置，与影像文件用同一个任务编号关联。

\paragraph{像控点与检查点的记录}每个点保存点号、现场照片、点位描述、平面坐标、高程、坐标参考系、测量设备、观测时段和质量等级。点位选在纹理稳定、边缘清楚、多个视角都能辨认的位置；坝顶栏杆、水面反光区和临时堆料区会随时间变化，不宜作控制点。检查点使用同样的记录格式，在空三求解阶段锁定为独立数据，最终验收时才参与误差统计。

控制点的布设按工程对象的几何层次考虑：坝轴线两端控制整体方向，坝顶和坝脚控制高差与坡面，闸室、溢洪道和廊道出入口控制局部构件，场景边界控制模型裁切和瓦片范围。每一层至少保留一组独立检查点，误差统计才能分别回答“整体是否平移”“局部是否变形”“边界是否翘曲”。所有点都布在坝顶时，平面指标可能很好，坝脚和峡谷侧壁的失真却发现不了。项目报告把点位分层、点数和每层的最大残差列成表。

\paragraph{内业检查}处理从影像完整性检查开始：核对文件数量、命名、时间戳、相机姿态、POS 轨迹和坐标单位，统计模糊、过曝、欠曝、遮挡和重叠不足的影像。检查结果分为“可进入空三”“需补拍或重采”“隔离待人工复核”三类，原始影像和检查报告只读保存。匀光匀色保持地物纹理的相对关系，处理参数和输出版本写入报告，增强操作不覆盖原始像素。

\paragraph{空三的四个阶段}第一阶段检测每幅影像的特征点，在相邻影像和交叉航带之间匹配连接点，并经几何一致性检验剔除错误对应。第二阶段读取相机内方位元素、镜头畸变参数和 POS 初值，建立观测方程；相机标定文件注明获取日期、适用镜头和坐标单位。第三阶段把像控点作为带权观测加入光束法平差，迭代估计相机外方位元素、连接点坐标和必要的系统改正；平差报告列出迭代次数、单位权中误差、控制点残差和未参与平差的检查点清单。第四阶段用独立检查点验证成果；误差集中在某条航带或某个高差突变区时，回查航带重叠、POS 时间同步、控制点坐标和镜头标定，而不是调高软件的容差。

空三参数的调整采用小步回归：冻结通过检查的影像、控制点和相机标定版本，每次只改变一个参数组（连接点匹配阈值、POS 权重或畸变模型），输出平差迭代、控制点残差和独立检查点的$m_p$、$m_h$，与上一版本比较。整体误差下降而某一航带误差上升时，保留两版结果并分析原因。重复飞行或季度更新时，稳定的控制点作为跨期公共点，新增检查点用来检验地表变化；发生过施工变化的点位标注“不可跨期比较”，避免把工程改造误判为摄影测量误差。

\paragraph{密集匹配与中间成果}密集匹配之前冻结通过空三验收的相机和点云版本。水面、玻璃和植被边缘纹理不足，会出现空洞；坝体混凝土纹理重复，会出现条纹；狭窄廊道视角不足，可能生成错误表面。处理报告记录点云密度、异常点剔除规则、空洞填补策略和人工修补区域；用于变形监测的几何不做视觉平滑。从点云到网格的每一步保留可回放的中间成果：点云阶段记录坐标基准、密度和分类规则，网格阶段记录三角形数量、孔洞填补和简化误差，纹理阶段记录影像选择、接缝处理和压缩参数，切片阶段记录层级、包围盒、瓦片格式和服务版本。坝面、闸门和监测仪器等关键对象同时保存原始高密度模型与发布用的轻量模型，用稳定的对象编码关联：平台运行时使用轻量瓦片，专业复核时回到原始模型核对尺寸和位置。

\paragraph{语义检查}把坝顶、坝脚、溢洪道边墙和库岸水线分别作为可查询对象，检查对象编码是否与 BIM 或 GIS 目录一致；同一构件在不同 LOD 里的几何简化，核对其中心线、端点和高程是否在允许范围内。监测仪器的位置与实测点号一一对应，模型更新只替换几何和纹理，不改变历史观测的对象主键。模型表面与控制点一致而对象属性错配时，成果退回语义校核环节。

\paragraph{验收结论}交付验收覆盖几何精度、坐标基准、对象语义、版本血缘和服务性能五个方面，分别保存检查数据、误差报告、属性映射表和加载测试结果。平面限差、高程限差、中误差计算口径和粗差处置规则作为一个整体签字确认；项目采用地方细化规范或合同中更严的要求时，注明版本、条款和适用范围；规范版本变化后重新计算检查点并保留旧版报告。验收单给出“通过、限期整改、退回重算”三种结论之一，连同责任人、完成日期和复核证据：通过表示精度和语义均满足适用条款；限期整改表示问题已经定位，且不影响隔离范围以外的成果；退回重算表示控制网、基准或关键对象存在系统性风险。平台的发布服务只读取“通过”版本。每次复核保存输入清单、软件环境和脚本版本，同一批检查点在另一台机器上应得到相同结果；跨期更新的报告注明哪些差异来自新采集、哪些来自处理参数。



\section{观测图表的更新预算、缓存与状态恢复}\label{app:ext-ch07-chart-runtime}

\paragraph{对应正文}第7章相关小节，具体入口见下文说明。

第7章7.2.7节说明了渲染器选择和更新节奏。图表从课堂演示走向长期运行的值班页面时，还会遇到下面几类问题。

\paragraph{缩放、查询与缓存}\texttt{dataZoom}只改变视图范围，浏览器里不必保留多年的原始数据。用户拖动到尚未缓存的时间段时，控制器按当前范围和合适的聚合粒度向服务端请求，并取消上一个未完成的请求；新应答返回后核对时间窗版本，晚到的旧应答不覆盖用户已经选定的新范围，做法与4.5.4节的序号核对相同。缓存键至少包含\texttt{assetId}、起止时间、聚合粒度、质量过滤条件、规则版本和时区，少了其中任何一项，同一条曲线在不同过滤条件下就可能被错误复用。在开发工具里显示缓存命中率、查询耗时和返回点数，可以看到“缩放流畅”主要取决于服务端按粒度聚合，而不是前端的某个技巧。

\paragraph{联动事件的一致性}图表点击、三维拾取和列表选择走同一个事件模型，事件里带着\texttt{assetId}、\texttt{occurredAt}、当前时间窗、质量码和来源组件。控制器先比较事件版本，再决定是否更新其他视图；给事件设幂等键，重复点击只更新一次焦点。指针移动产生的高频预览可以节流，确认预警、提交工单和导出报告不能节流。时间窗输入采用防抖并显示加载状态。服务端返回时如果规则版本已经变化，界面提示“按新规则重新计算”，不要静默替换颜色。地图、图表与三维场景之间共用一个选择状态对象，例如\texttt{\{assetId, timeRange, qualityFilter, scenarioVersion\}}；切换场景版本时清空或重新校验它，避免把上一版本的测点高亮到新模型上。

\paragraph{保护用户的上下文}数据刷新不重置当前的缩放范围、选中测点、键盘焦点和已展开的详情。规则版本改变时先保留旧视图并提示差异，用户确认后再切换。导出的图片或表格使用与屏幕相同的时间窗和过滤条件，并把采样周期、聚合方法、质量过滤和规则版本写入元数据，值班员看到的屏幕、分析员导出的报告和事后回放的记录才能互相核对。

\paragraph{刷新与断线后的恢复}浏览器刷新、网络短暂中断和标签页挂起都会打断实时页面。状态层可以保存一份快照：最近一次通过校验的时间窗、各序列的摘要、规则版本和最后一个事件游标。恢复时先显示“数据恢复中”，再按游标向服务端补拉增量，补拉的数据重新经过去重、迟到判断和质量规则。快照里不保存整个 ECharts 配置对象，也不保存临时的颜色，否则旧的合并结果会被当成新状态。补拉完成之前，断线期间在图上保持缺测阴影，不画成连续曲线。值班员确认预警的操作由后端记录确认人、时间和规则版本，刷新页面或换一台设备后确认状态仍在。

\paragraph{可以复用的回归数据}图表部分准备五组数据：单点连续到达；同一事件重复到达；乱序与迟到；含\texttt{missing}区间；一条序列删除后再加入。逐组检查曲线是否按\texttt{id}更新、旧图例是否清理、缺测是否断线、参考线是否与参数版本一致、联动能否在图表与三维之间往返。再在窄屏、深色主题、灰度打印和纯键盘操作下各测一遍，并确认实例销毁后没有残留的订阅。地图部分另加边界外的点、跨瓦片的点和坐标转换失败的点。每个失败用例保存输入快照、预期图形和实际图形，升级 ECharts 版本时用同一组数据回归。三维交互的回归再覆盖原点变更、LOD 切换和时间回放三者的组合，并记录相机位置、视场角、设备像素比、窗口尺寸、渲染器和浏览器版本，同一个点击坐标在不同设备上的结果才能对比；命中的三维对象、图表数据下标和事件时间三者不一致时，先修坐标或状态契约，不要靠调大拾取半径掩盖。



\section{颜色令牌、验收矩阵与色板版本}\label{app:ext-ch07-color-tokens}

\paragraph{对应正文}第7章相关小节，具体入口见下文说明。

第7章7.2.9节给出了色板类型、对比度公式和冗余通道的做法。多人协作、多主题的平台还需要把颜色当作配置来管理。

\paragraph{颜色令牌}先建立颜色令牌表，再让组件引用令牌名称，页面里不直接写十六进制色值。令牌分六组：背景、正文、边框、顺序型色阶、预警等级和质量状态。背景令牌有浅色和深色两套取值，正文令牌同时记录普通文字和大号文字的对比度结果。切换大屏主题时只替换一组令牌并重新计算对比度，不会出现地图已经变暗、图例仍是浅色文字的局部失配。每个令牌还记录适用对象和禁用对象：红色预警令牌只表示需要立即处置的业务等级，不兼作“数据无效”；蓝色预警令牌不兼作链接色或选中色。同一个点既是\texttt{suspect}又是橙色预警时，详情里质量码、等级、规则版本和处置建议分列为独立字段，颜色只是辅助。

\paragraph{验收矩阵}一张静态截图通过对比度检查，不能说明运行时也可读。按“背景×状态×组件×通道”建立矩阵：背景取浅色、深色、灰度和投影；状态取 NONE、BLUE、YELLOW、ORANGE、RED 与\texttt{valid}、\texttt{suspect}、\texttt{missing}的组合；组件取地图点、折线、柱状图、表格、按钮和弹窗；通道检查颜色、文字、图标、纹理和键盘焦点。每一格记录前景色、背景色、对比度、有无非颜色提示、读屏软件读出的文本和截图编号。某一格只依赖颜色时记为不通过，补上文字或形状后再复查。

\paragraph{状态变化时的可读性}预警等级从 NONE 升为 BLUE，界面除了变色，还显示“出现需关注变化”的文字，并在时间轴上留下状态变化的记录；从 BLUE 升为 YELLOW，同时更新处置入口和确认状态。质量码从\texttt{valid}变为\texttt{missing}时折线断开；恢复为\texttt{valid}后缺测区间保持空白，详情里显示恢复时间，补传的数据不伪装成连续的实时数据。动画时间要短并且可以暂停，闪烁只用于需要立即确认的事件。

\paragraph{色板版本}色板与预警规则版本一起放进配置仓库。每次调整蓝、黄、橙、红的色值，记录变更原因、影响的页面、对比度复测结果和导出报告样例；历史报告按生成时的色板版本解释，同一个等级在不同月份的报告里含义一致。前端加载令牌时发现版本不匹配，开发环境给出明显提示，生产环境回退到最近一次通过验收的版本并上报监控。

\paragraph{打印与导出}黑白打印时用线型、填充纹理和缩写文字区分等级。导出的表格把颜色还原成“预警等级”和“质量码”两列，并在表头给出解释。图片的替代文本按“对象—数值—质量—预警—时间”的顺序生成。验收时对比屏幕、打印件和导出文件，检查等级、质量与时间在三处是否能被一致地识别。



\section{第8章数据库与部署的运行维护练习}\label{sec:appc-ch08-ops}

\paragraph{对应正文}第8章相关小节，具体入口见下文说明。

第8章核心路线只用到建表、幂等写入、条件更新和按事件追溯。本节收录其余的数据库练习和部署运维要点。表结构见第8章8.3.3节；所有语句都在隔离的练习库里执行，保留和压缩策略会删数据、改存储格式，不要对着有真实数据的库试。

\subsection{迁移版本、清单查询与空间查询}

迁移脚本要能重复执行，执行过的版本记在数据库里。清单\ref{lst:ch08-sql-migration}用一张轻量的版本表记迁移名称和校验摘要；发布工具另记执行人、开始时间、结束时间和失败日志。迁移中途失败时，恢复脚本只回滚本次版本，不删已有的业务数据。

\begin{lstlisting}[label={lst:ch08-sql-migration},language=SQL,caption={可重复执行的迁移版本记录}]
CREATE TABLE IF NOT EXISTS schema_migration (
    version      text PRIMARY KEY,
    description  text NOT NULL,
    checksum     text NOT NULL,
    applied_at   timestamptz NOT NULL DEFAULT now(),
    applied_by   text NOT NULL
);

INSERT INTO schema_migration(version, description, checksum, applied_by)
-- checksum 由发布流水线对迁移脚本正文计算后填入，禁止手工编造
VALUES ('QY-DDL-001', 'core asset reading warning ddl',
        'sha256:9f2d1c0a...(流水线填入)', current_user)
ON CONFLICT (version) DO NOTHING;
\end{lstlisting}

对象清单是最频繁的读路径之一。清单\ref{lst:ch08-sql-asset-list}用状态和类型过滤有效测点，只取前端需要的列：不写\texttt{SELECT *}，以后加了列接口载荷也不会悄悄变大，数据库也更容易用上覆盖索引。分页用稳定的\texttt{asset\_id}游标（\ref{app:ext-perf}节），页码在数据增删之后会漂移。

\begin{lstlisting}[label={lst:ch08-sql-asset-list},language=SQL,caption={按类型和状态读取测点清单}]
SELECT asset_id, display_name, asset_type, unit,
       ST_X(geometry) AS longitude,
       ST_Y(geometry) AS latitude,
       elevation_m
FROM asset
WHERE active = true
  AND asset_type = :asset_type
  AND (:after_id IS NULL OR asset_id > :after_id)
ORDER BY asset_id
LIMIT :page_size;
\end{lstlisting}

空间查询先说清输入多边形的坐标系。清单\ref{lst:ch08-sql-spatial-query}把外部边界转换到 4490，先用 GiST 索引支持的包围盒过滤（\texttt{\&\&}），再做精确的\texttt{ST\_Within}判断。输入无效时接口在参数校验阶段就返回错误；交给数据库执行，得到的是空结果，会被当成“区域内没有测点”。

\begin{lstlisting}[label={lst:ch08-sql-spatial-query},language=SQL,caption={影响区内测点的空间查询}]
WITH area AS (
    SELECT ST_Transform(
        ST_SetSRID(ST_GeomFromText(:wkt), :input_srid), 4490
    ) AS boundary
)
SELECT a.asset_id, a.display_name, a.asset_type
FROM asset AS a CROSS JOIN area
WHERE a.active
  AND a.geometry && area.boundary
  AND ST_Within(a.geometry, area.boundary)
ORDER BY a.asset_id;
\end{lstlisting}

工单到期查询用数据库的当前时间和明确的状态集合。清单\ref{lst:ch08-sql-due-work-order}排除已取消或已完成的工单，连预警等级一起返回，值班员先处理等级高、临近截止的。更新状态时带版本号（8.3.4节的条件更新），两个角色才不会互相覆盖。

\begin{lstlisting}[label={lst:ch08-sql-due-work-order},language=SQL,caption={即将到期工单查询}]
SELECT w.work_order_id, w.warning_id, w.owner_role,
       w.due_at, w.action, p.level
FROM work_order AS w
JOIN warning AS p ON p.warning_id = w.warning_id
WHERE w.status = 'in_progress'
  AND w.due_at <= now() + INTERVAL '2 hours'
ORDER BY p.level DESC, w.due_at;
\end{lstlisting}

模型运行审计把版本和输入快照一起返回。清单\ref{lst:ch08-sql-model-audit}按模型名称和时间窗筛出成功与失败的运行；失败的运行只用来定位故障，下游不拿它当结果。结果要重新发布时新建一条运行记录，旧结果不覆盖。

\begin{lstlisting}[label={lst:ch08-sql-model-audit},language=SQL,caption={模型运行审计查询}]
SELECT run_id, model_name, model_version, status,
       started_at, finished_at, input_snapshot
FROM model_run
WHERE model_name = :model_name
  AND started_at >= :start_at
  AND started_at < :end_at
ORDER BY started_at DESC;
\end{lstlisting}

建表顺序见第8章8.3.3节。已经产生证据的表不能靠删表来回滚：先停止写入，导出审计快照，再把迁移标记为失效。升级字段按“增加列—双写—回填—切换—观察—清理”六步走，观察期内保留旧列，直到备份和一次恢复演练都成功。

\subsection{连续聚合、保留与压缩}

连续聚合是TimescaleDB提供的一种物化视图：数据库按时间桶预先算好平均值、最值和样本数，并随新数据自动刷新，重复读取趋势时不必每次扫描原始观测。清单\ref{lst:ch08-sql-cagg}在\texttt{reading}上建立十五分钟聚合\texttt{reading\_15m}，只纳入\texttt{valid}观测，并保留样本数，页面可以据此解释“平均值为什么没有算上某条数据”。刷新策略要有一个回看窗口来覆盖迟到的消息；例中回看30天、延迟5分钟是教学参数，实际窗口应覆盖通信补传和消息重投的最长时间。

\begin{lstlisting}[label={lst:ch08-sql-cagg},language=SQL,caption={十五分钟连续聚合视图与刷新策略}]
CREATE MATERIALIZED VIEW reading_15m
WITH (timescaledb.continuous) AS
SELECT asset_id,
       time_bucket('15 minutes', occurred_at) AS bucket,
       avg(value) AS value_avg, min(value) AS value_min,
       max(value) AS value_max, count(*) AS sample_count
FROM reading
WHERE quality = 'valid'
GROUP BY asset_id, bucket
WITH NO DATA;

SELECT add_continuous_aggregate_policy(
    'reading_15m',
    start_offset => INTERVAL '30 days',
    end_offset => INTERVAL '5 minutes',
    schedule_interval => INTERVAL '5 minutes'
);
\end{lstlisting}

聚合结果适合重复读取的统计，当不了证据。预警计算可以读十五分钟平均值来降低延迟，触发处置之前保留对应的原始时间窗、样本数和质量分布。聚合表不手工编辑；数值要修正，回到原始观测追加订正版本，再刷新聚合。

连续聚合查询只读已物化的时间桶。清单\ref{lst:ch08-sql-aggregate-query}把桶边界、样本数一起返回，前端在样本不足时显示“数据稀疏”，不画一条看似连续的折线。数据库统一用 UTC 存时间，展示层按用户时区转换，跨时区部署才不会乱。

\begin{lstlisting}[label={lst:ch08-sql-aggregate-query},language=SQL,caption={读取十五分钟连续聚合结果}]
SELECT bucket, value_avg, value_min, value_max, sample_count
FROM reading_15m
WHERE asset_id = :asset_id
  AND bucket >= :start_at
  AND bucket < :end_at
ORDER BY bucket;
\end{lstlisting}

迟到的消息会让最近的聚合桶需要重算。清单\ref{lst:ch08-sql-refresh-aggregate}是手工刷新某个窗口的语句，运维员确认补传范围后执行，变更记录里写明起止时间。日常刷新仍由自动策略负责；手工刷新是补救手段，不是绕过质量审核的常规入口。

\begin{lstlisting}[label={lst:ch08-sql-refresh-aggregate},language=SQL,caption={补传后的连续聚合刷新}]
CALL refresh_continuous_aggregate(
    'reading_15m', :refresh_start, :refresh_end
);

INSERT INTO schema_migration(version, description, checksum, applied_by)
VALUES (:audit_version, 'late data aggregate refresh', :checksum, current_user);
\end{lstlisting}

冷数据的保留期服从证据保留期限。清单\ref{lst:ch08-sql-retention}用 TimescaleDB 的策略对象表达自动清理窗口；730 天是教学取值，生产项目的值从法规、合同和审计要求里来，删除前先完成备份校验。某段数据还被未关闭的工单引用时，业务服务先延长它的保留期。

\begin{lstlisting}[label={lst:ch08-sql-retention},language=SQL,caption={时序数据保留策略示例}]
SELECT add_retention_policy(
    'reading', drop_after => INTERVAL '730 days'
);

-- 发布前由项目配置替换保留窗口，并核对审计与备份要求
SELECT hypertable_name, num_chunks, compression_enabled
FROM timescaledb_information.hypertables
WHERE hypertable_name = 'reading';
\end{lstlisting}

压缩避开还在频繁更新的热数据。清单\ref{lst:ch08-sql-compression}按时间列组织压缩，30 天以前的块才压；质量修订窗口没关闭的数据不压，压缩块上的更新代价很高。压缩不是删除，恢复演练照样要验证索引、连续聚合和原始版本都能读。

\begin{lstlisting}[label={lst:ch08-sql-compression},language=SQL,caption={按时间块配置时序压缩}]
ALTER TABLE reading SET (
    timescaledb.compress,
    timescaledb.compress_segmentby = 'asset_id',
    timescaledb.compress_orderby = 'occurred_at DESC'
);

SELECT add_compression_policy(
    'reading', compress_after => INTERVAL '30 days'
);
\end{lstlisting}

时间分块的意义是把“最近的数据最常用”变成物理上的局部性：查最近一小时不碰几年前的块，归档两年前的数据不锁住正在写入的块。块的大小用有代表性的写入数据压测后确定：太小，规划器要管理大量对象；太大，压缩和备份的粒度变粗。

\subsection{健康检查、备份登记、权限视图与修订审计}

数据库健康检查找的是约束被违反的前兆。清单\ref{lst:ch08-sql-health-check}统计缺坐标、质量码不在集合内和时间戳在未来的记录；它不改数据，只出报告。每日检查结果和发布版本关联起来，出现异常时才分得清是设备问题、迁移问题还是程序问题。

\begin{lstlisting}[label={lst:ch08-sql-health-check},language=SQL,caption={核心表健康检查}]
SELECT 'asset_without_geometry' AS check_name, count(*) AS failures
FROM asset WHERE geometry IS NULL
UNION ALL
SELECT 'reading_bad_quality', count(*)
FROM reading WHERE quality NOT IN ('valid','suspect','missing')
UNION ALL
SELECT 'reading_in_future', count(*)
FROM reading WHERE occurred_at > now();
\end{lstlisting}

备份恢复演练验证的是“业务链能不能恢复”，不是“文件在不在”。清单\ref{lst:ch08-sql-backup-scope}只登记快照的元信息，导出数据和权限由备份工具负责。演练完成后用一条事件 ID 串起测点、观测、预警、工单和模型运行，确认外键和 JSONB 证据都能读。

\begin{lstlisting}[label={lst:ch08-sql-backup-scope},language=SQL,caption={备份快照登记}]
CREATE TABLE IF NOT EXISTS backup_snapshot (
    snapshot_id  text PRIMARY KEY,
    started_at   timestamptz NOT NULL,
    finished_at  timestamptz,
    object_count integer NOT NULL,
    checksum     text NOT NULL,
    verified     boolean NOT NULL DEFAULT false
);
INSERT INTO backup_snapshot(snapshot_id, started_at, object_count, checksum)
VALUES (:snapshot_id, now(), :object_count, :checksum);
\end{lstlisting}

权限查询只返回当前角色可见的对象。清单\ref{lst:ch08-sql-role-view}用视图把敏感字段包起来，报表用户看得到编码、等级和统计结果，看不到用不着的个人信息和设备密钥。视图是数据库这一层的最小权限，应用层授权和数据库审计照样要有。

\begin{lstlisting}[label={lst:ch08-sql-role-view},language=SQL,caption={面向报表角色的最小视图}]
CREATE OR REPLACE VIEW warning_report AS
SELECT warning_id, asset_id, level, evaluable,
       score, reason, status, created_at
FROM warning
WHERE status <> 'closed' OR created_at >= now() - INTERVAL '90 days';

GRANT SELECT ON warning_report TO report_reader;
\end{lstlisting}

异常数据的修复要能审计。清单\ref{lst:ch08-sql-revision-audit}在写入新观测版本之前先登记修订前后的值、原因和批准人；没有批准号的修订，接口拒绝提交。审计表可以单独归档，前端只显示最新值不是删它的理由。

\begin{lstlisting}[label={lst:ch08-sql-revision-audit},language=SQL,caption={观测修订审计登记}]
CREATE TABLE IF NOT EXISTS reading_revision_audit (
    audit_id       bigserial PRIMARY KEY,
    event_id       text NOT NULL,
    old_value      numeric,
    new_value      numeric,
    old_quality    text NOT NULL,
    new_quality    text NOT NULL,
    reason         text NOT NULL,
    approved_by    text NOT NULL,
    audited_at     timestamptz NOT NULL DEFAULT now()
);
\end{lstlisting}

数据库的观察指标和业务指标对着看。清单\ref{lst:ch08-sql-observability}从系统视图读超表大小、索引大小和最近一条数据的时间，给压测和运维看板提供证据，不直接下“系统健康”的结论。指标异常时同时看 Kafka 积压、质量码分布和工单延迟，只看数据库 CPU 看不出问题在哪。

\begin{lstlisting}[label={lst:ch08-sql-observability},language=SQL,caption={数据库容量与新鲜度观察}]
SELECT hypertable_name, table_bytes, index_bytes,
       total_bytes
FROM hypertable_detailed_size('reading');

SELECT max(occurred_at) AS latest_reading,
       now() - max(occurred_at) AS data_lag
FROM reading;
\end{lstlisting}

备份同时覆盖结构、数据、权限和扩展。PostGIS 的空间列、TimescaleDB 的超表元数据、连续聚合策略和角色授权，缺任何一项恢复后都可能查不出来。全量备份之外保留增量或归档日志，定期在隔离环境恢复一次。恢复后按事件标识查测点、原始观测、质量结论、预警记录、工单回执和模型输入，再检查角色授权对不对、新的时间块能不能继续写入；有一项失败就记为恢复缺口。

数据库权限按最小必要分配：采集角色只写观测和发件箱；质量角色能改质量标记，不能改原始值；预警角色能创建和关闭预警；值班员能确认工单并提交回执；专业分析员能读脱敏的历史数据；运维员才能做分区、压缩和恢复。数据库角色与应用角色一一对应，凭据放在密钥服务里，代码和教材示例不出现真实口令。审计日志记角色、对象、操作、时间、请求追踪号和结果摘要，数据库行为才能和 Kafka 消息、HTTP 请求、前端操作对上。

模式演进先考虑兼容。加可空列通常向后兼容；删列和改枚举要经过观察期；大表加非空约束先回填并验证，再分阶段切换。实体类和 DTO 不另起一套字段名，数据库字典是唯一来源，代码审查时逐列核对。后端把\texttt{asset}映射成 Java 实体时保留\texttt{asset\_id}、\texttt{geometry}和\texttt{metadata}的原语义，空间类型的转换策略写明；S3 生成示例数据时按同一套编码规则和质量码集合生成。

测试覆盖约束、查询计划和故障恢复三层。单元测试验证质量优先级和分类结果；集成测试验证外键、唯一键、空间坐标系、超表写入和连续聚合；端到端测试验证从事件接入到工单回执的链路。性能测试用与目标测点规模相近的数据分布，分别记高峰写入、最近查询、空间查询、预警队列和聚合刷新的延迟。故障测试包括消息重复、乱序、数据库短暂不可用、超表块压缩失败和恢复后重放，每个故障写明预期的质量码、状态转移和用户看到的提示。

容量规划从测点数量、采样周期和保留期限推：先算每天的事件行数，再加索引、JSONB 证据、备份副本和压缩后的增长系数；写入峰值考虑设备同时上线、补传和消息重投。算出来的只是初始值，最终用接近生产分布的回放数据压测。压测记平均吞吐之外的 P95、P99 延迟、锁等待、超表块数量、连续聚合刷新滞后，以及故障重试后积压恢复的时间。

\subsection{告警规则、容量与备份恢复演练}

健康检查的三个层次见第8章8.6.1节；业务烟测在发布后和每日交接时各跑一次。清单\ref{lst:ch08-alert-rules}是一组告警规则的基线。

\begin{lstlisting}[language=YAML,caption={监控指标与告警规则基线},label={lst:ch08-alert-rules}]
groups:
  - name: reservoir-platform
    interval: 30s
    rules:
      - alert: ApiReadinessFailed
        expr: api_readiness == 0
        for: 2m
        labels: {severity: critical, owner: ops}
      - alert: ReadingIngestLag
        expr: reservoir_reading_ingest_lag_seconds > 300
        for: 5m
        labels: {severity: warning, owner: duty}
      - alert: TimescaleDiskPressure
        expr: reservoir_db_free_bytes < 20000000000
        for: 10m
        labels: {severity: critical, owner: ops}
      - alert: WarningAckTimeout
        expr: reservoir_warning_unacked_seconds > 900
        for: 5m
        labels: {severity: warning, owner: duty}
\end{lstlisting}

清单里的规则按基础设施（\texttt{ApiReadinessFailed}、\texttt{TimescaleDiskPressure}）、数据链路（\texttt{ReadingIngestLag}）和业务闭环（\texttt{WarningAckTimeout}）分组，每条告警有责任角色（\texttt{owner}）和处置手册。阈值结合补传的最大重投周期、TimescaleDB 压缩窗口和工单确认时限来定，再用历史分位数校准。告警和恢复都记时间、规则版本、当前值和证据链接；同一根因的指标按服务和时间窗聚合，否则一次故障会刷出几十条告警。

容量评估覆盖 API 并发、Kafka 分区、Redis 内存、数据库写入吞吐和时序增长。\cpStationCount 个测点的教学规模单节点就能跑，推不出生产容量；压测把补传、模型运行和备份的并发都算进去。扩容前后用同一场景复测，记消息最老年龄、WAL 增长、磁盘剩余比例和查询 P95，看容量提升有没有换来数据延迟或权限错误。

备份策略回答三个问题：可接受的数据丢失窗口（RPO）是多少，一致性点在哪里，怎样验证恢复。PostgreSQL 用全量加 WAL 归档（WAL 是 PostgreSQL 的预写日志，记录每一次修改，配合全量备份可以恢复到任意时刻）；TimescaleDB 的压缩和分区任务与备份窗口错开；对象存储保存模型、纹理、报告和备份清单的版本；Kafka 保留足够长的日志供重放，重放靠\texttt{event\_id}唯一约束去重。Redis 只放可重建的缓存，不进备份范围。清单\ref{lst:ch08-backup-rehearsal}把这些写成一份演练配置。

\begin{lstlisting}[language=YAML,caption={备份恢复演练与健康检查清单},label={lst:ch08-backup-rehearsal}]
backup: {schedule: "15 2 * * *", retention_days: 35, mode: "base-plus-wal", verify: [checksum, restore-schema, latest-reading-time]}
object_store: {versioning: true, retention_days: 90}
kafka: {topic: reservoir.reading.v1, replay_window_hours: 72}
rehearsal:
  target: isolated-restore
  steps: [stop-writes, restore-postgres-and-wal, restore-object-manifest, start-api-readonly, verify-trace, replay-with-event-id-deduplication, run-smoke-test]
  pass_conditions: [latest_reading_within_rpo, api_readiness_up, no_duplicate_warning, audit_trace_continuous]
\end{lstlisting}

清单\ref{lst:ch08-backup-rehearsal}把“有备份”变成可观察的证据。开始前冻结变更窗口，记下最后一个观测、预警和工单 ID；恢复后先以只读模式校验对象关系、时序最新时间和审计链，再开放写入。Kafka 只重放恢复点之后、数据库里没有的事件；缓存由 API 按数据库重建。\texttt{pass\_conditions}里任何一条不满足，演练就是不通过，记录里写明是缺备份、迁移出错还是有解释不了的事件。

生产故障按观测、隔离、恢复、复盘的顺序处置。运维员先看探针和指标是不是同时异常，一个浏览器出问题不等于全局故障；再按运行手册摘掉异常实例或暂停消费，保住数据库和消息证据；恢复后用业务烟测验证从测点到工单的链路，专业分析员复核质量码和预警有没有重复计算。每次演练记下恢复时长、数据丢失窗口、告警命中率和人工确认点，反过来校准保留策略、探针阈值和容量预算。

环境一致性还要查时区、字符集、容器用户和文件权限。数据库统一用 UTC 存发生时间，展示层按值班员时区转换；所有容器以非 root 用户运行，备份目录只给备份进程写权限。发布前用同一份配置渲染开发、测试、生产三套文件，比较键的集合，多了或少了变量就停止发布。配置差异写进决策记录，不靠运维员记着。

服务启动分迁移、就绪、接流量三个阶段。迁移在独立的一次性容器里执行并记录版本，API 进程启动时只验证当前版本，多个副本不同时改表结构。就绪探针恢复后逐步放量：先让一个实例接只读查询，再开观测写入和工单写入。错误率或延迟超过发布阈值就停止放量、保留旧副本，回不回滚由值班员和审批人一起定。

监控标签只用稳定的维度：测点编号、工程编号、服务名和环境。用户姓名、完整请求参数和原始观测值不进指标，一是基数太高，二是泄露敏感信息。日志记请求号、状态码、耗时和错误类型，详细堆栈只进受控的日志库。运维看板同时展示基础设施、数据质量和业务闭环，值班员看到的是能据此行动的提示，不是解释不了的技术计数。

告警处理有确认、抑制、升级、恢复四个状态。值班员确认不等于故障解决了，烟测和业务链路验证通过才能关闭事件；维护窗口内的预期告警可以抑制，抑制规则有开始和结束时间。关键告警超过确认时限自动升级给运维员，超过处置时限再通知审批人。每次升级保留原始告警和责任转移记录，交接后上下文才不丢。

Kafka 重放和定时补传都可能造成重复事件。服务端靠\texttt{event\_id}唯一约束和事务保证幂等；演练时运维员主动重放一条已经消费过的消息，看会不会重复创建预警或工单。重放窗口覆盖消息的最大重投周期，分区、偏移量和过滤条件写进演练记录。发现重复时先暂停消费者，保留现场，由开发人员依据审计链修复，不直接删数据库记录。

备份保留多久，由恢复目标、法规要求和容量预算一起定。太短，长周期趋势无法复盘；太长，成本和恢复扫描时间上去。清理任务只删已通过校验且超出保留期的版本，删之前生成清单等审批。压缩、分区和归档策略每次变更先在恢复环境验证，线上压缩不能拖慢实时写入。

部署验收包含故障注入，不只走正常路径：停掉 Redis，看 API 是否降级为数据库查询；暂停 Kafka 消费，看积压告警和恢复后的幂等；限制数据库磁盘，看容量告警和写入保护；撤销一个角色，看 403 是否正确显示。每项注入设最长持续时间和回滚动作，演练完确认没有遗留进程、临时密钥或没关的维护窗口。

部署实验提交这些记录：编排文件的版本、渲染后的非敏感配置、探针应答、告警触发与恢复时间、备份校验摘要、恢复后的业务烟测和变更审批号。教师据此判断部署可不可重复、数据可不可恢复、权限有没有效、运维责任清不清楚。迁到多节点编排环境时，再补节点故障、持久卷调度和跨节点网络的恢复测试。


\section{教学组织与学时建议}\label{online:teaching-guide}

\paragraph{对应正文}前言的学习顺序与阶段表；以下安排供教师按教学条件选用。

三个学时方案按学生最终能做到什么来编制，学时少的方案少走几个阶段，已走过的阶段仍然完整。方案之间的差别在于沿表\ref{tab:preface-version-stage}走到哪个阶段：32学时走到 S3 中点，48学时走到 S5，56学时走到 S6 并完成部署验收。每个方案都包含表\ref{tab:preface-core-optional}左列的对应部分，右列的选学专题留给课程设计或自学。

\begin{itemize}
\item \textbf{32学时方案（到 S3 中点）：}学生能解释一次请求的全过程，把一句含糊需求写成验收条件，在教学接口上做出带四种状态的监测页面，写出第一个 HTTP 接口并把它接到数据库。不讲三维场景、曲线联动和预警闭环；第6--8章可作为课程设计的阅读材料。表\ref{tab:preface-32-hours}是章节分配，适合专业方向课或选修课。
\begin{table}[htbp]
\centering
\caption{32学时方案章节分配（到 S3 中点）}
\label{tab:preface-32-hours}
\begin{tabular}{p{2.8cm}p{1.4cm}p{8.2cm}}
\toprule
章节或环节 & 学时 & 课堂重点 \\
\midrule
第1章 & 2 & 平台组成、三类规范、一次查询的六个环节（S0）、预警闭环 \\
第2章 & 4 & 用例与异常流、数据字典、从含糊需求到验收条件 \\
第3章 & 4 & 需求分给模块、接口契约、走查一次查询 \\
第4章 4.1--4.5 & 9 & 静态页面（S1）、请求与四种页面状态、竞态故障（S2 核心） \\
第5章 5.1--5.4 & 11 & 第一个接口（S3 起点）、数据库查询与三层（S3 中点） \\
复习与个人答辩 & 2 & 对着自己的页面和接口解释一次失败 \\
\bottomrule
\end{tabular}
\end{table}

\item \textbf{48学时方案（到 S5）：}在32学时的基础上完成认证与授权（S3 终点）、坝体场景与测点绑定（S4）、曲线与三维联动（S5）。第8章只读8.1--8.2，了解业务链怎样由已学模块接成，不做预警与工单的实现。表\ref{tab:preface-48-hours}给出分配。
\begin{table}[htbp]
\centering
\caption{48学时方案章节分配（到 S5）}
\label{tab:preface-48-hours}
\begin{tabular}{p{2.8cm}p{1.4cm}p{8.2cm}}
\toprule
章节或环节 & 学时 & 课堂重点 \\
\midrule
第1章 & 2 & 平台组成、三类规范、一次查询的六个环节（S0）、预警闭环 \\
第2章 & 4 & 用例与异常流、数据字典、从含糊需求到验收条件 \\
第3章 & 3 & 需求分给模块、接口契约、走查两条路径 \\
第4章 & 10 & S1 静态页面、S2 四种状态与竞态、Vue 路由与登录守卫 \\
第5章 5.1--5.6 & 13 & S3 起点、中点与终点：接口、数据库、三层、认证与授权 \\
第6章 6.1--6.2 & 6 & 几何体坝体、模型加载、坐标与高程、测点绑定（S4） \\
第7章 7.1--7.4 & 6 & 观测记录与质量码、第一条曲线、双向联动（S5） \\
第8章 8.1--8.2 & 2 & 读懂 PZ-07 业务链用到哪些已学模块 \\
复习与项目答辩 & 2 & 演示 S5 页面，解释一次切换测点时的请求过程 \\
\bottomrule
\end{tabular}
\end{table}

\item \textbf{56学时方案（到 S6 与部署）：}讲授第1--8章的核心与指导实践小节，第4、5章各占14学时，其中一半用于随堂实验；第8章完成质量门禁、预警定级、确认与工单闭环，并用 Compose 与冒烟测试做一次部署验收。第6章6.3--6.4、第8章8.5和第9章9.2.1--9.2.2作为选读。表\ref{tab:preface-56-hours}给出分配，讲授与实验之比约为30:26。
\begin{table}[htbp]
\centering
\caption{56学时方案章节分配（到 S6 与部署）}
\label{tab:preface-56-hours}
\begin{tabular}{p{2.8cm}p{1.2cm}p{1.2cm}p{7.0cm}}
\toprule
章节或环节 & 讲授 & 实验 & 课堂重点 \\
\midrule
第1章 & 2 & 0 & 一次请求的完整过程、平台分层、三类规范与业务闭环 \\
第2章 & 3 & 2 & 用例、异常流、数据字典与可验证需求 \\
第3章 & 3 & 2 & 需求分给模块、接口契约、两条走查记录、一份 ADR（架构决策记录） \\
第4章 & 7 & 7 & 页面、数据请求与四种状态、Vue 组件、路由与共享状态 \\
第5章 & 7 & 7 & 第一个 HTTP 接口、数据库与三层、校验与事务、认证与权限 \\
第6章 & 3 & 2 & 几何体坝体、坐标与高程、模型加载与测点绑定 \\
第7章 & 3 & 2 & 观测记录与质量码、第一条曲线、测点状态与联动 \\
第8章 & 2 & 3 & 质量门禁、四级定级、确认与工单闭环、部署冒烟测试 \\
第9章与答辩 & 0 & 1 & 按课程成果回顾、现场小修改与解释 \\
\bottomrule
\end{tabular}
\end{table}

\item \textbf{课程设计方案：}以第8章贯穿案例为中心，要求学生分组完成需求说明、架构设计、前端页面、后端接口和三维场景集成，可与毕业设计或综合实训衔接。
\end{itemize}


% BEGIN learning-layers
\section{各节学习层次}\label{online:layers}

\paragraph{说明}本表供教师安排课时使用。“核心”是课堂讲授与闭卷考核的范围；“指导实践”在工程骨架上完成，主要安排在随堂实验；“拓展”供课程设计与自学。只写一个层次的节，表示该节所有小节都属于这一层次。

\begin{longtable}{p{1.2cm}p{4.6cm}p{8.0cm}}
\toprule
节 & 标题 & 各小节层次 \\
\midrule
1.1 & 智慧水利的概念与发展 & 核心：1.1.1、1.1.2、1.1.3、1.1.5；拓展：1.1.4 \\
1.2 & 智慧水利平台架构体系 & 核心：1.2.1、1.2.2、1.2.3、1.2.4、1.2.5 \\
1.3 & 为什么水利平台需要软件工程 & 核心 \\
2.1 & 软件生命周期 & 核心 \\
2.2 & 软件生命周期模型与选型 & 核心：2.2.1、2.2.2、2.2.3、2.2.4；指导实践：2.2.5；拓展：2.2.6 \\
2.3 & 需求获取与理解 & 核心：2.3.1、2.3.2、2.3.3、2.3.4、2.3.5 \\
2.4 & 需求分析的任务与原则 & 核心：2.4.1、2.4.2 \\
2.5 & 结构化分析 & 核心：2.5.1、2.5.2、2.5.3、2.5.4；指导实践：2.5.5、2.5.6、2.5.7 \\
2.6 & 软件需求规格说明与验收 & 核心：2.6.1、2.6.4；指导实践：2.6.2；拓展：2.6.3 \\
3.1 & 软件架构及其描述 & 核心：3.1.1、3.1.2、3.1.3、3.1.4、3.1.5 \\
3.2 & 架构类型与分层 & 核心：3.2.1、3.2.2；指导实践：3.2.3 \\
3.3 & 接口与数据流 & 核心：3.3.1、3.3.2、3.3.3；指导实践：3.3.4 \\
3.4 & 模块化与依赖 & 核心：3.4.1、3.4.2、3.4.3；拓展：3.4.4 \\
3.5 & 架构选择 & 核心：3.5.1、3.5.2；拓展：3.5.3、3.5.4、3.5.5 \\
3.6 & 架构评审 & 核心：3.6.1；指导实践：3.6.2；拓展：3.6.3、3.6.4、3.6.5 \\
4.1 & Web基础概念与浏览器架构 & 核心 \\
4.2 & HTML基础与语义化 & 核心：4.2.1、4.2.2；指导实践：4.2.3 \\
4.3 & CSS样式与响应式布局 & 核心：4.3.1、4.3.2、4.3.3、4.3.4、4.3.5；指导实践：4.3.6；拓展：4.3.7、4.3.8 \\
4.4 & JavaScript基础编程 & 核心：4.4.1、4.4.2、4.4.3、4.4.4、4.4.5、4.4.7、4.4.8；指导实践：4.4.6 \\
4.5 & 数据请求与页面状态 & 核心：4.5.1、4.5.2、4.5.3、4.5.4、4.5.10；指导实践：4.5.5、4.5.6、4.5.7；拓展：4.5.8、4.5.9 \\
4.6 & Vue 3组件化开发 & 核心：4.6.1、4.6.2；指导实践：4.6.3、4.6.4 \\
4.7 & Vue Router与Pinia应用组织 & 指导实践：4.7.1、4.7.2、4.7.3、4.7.4 \\
4.8 & 调试、构建与阶段验收 & 核心：4.8.1；指导实践：4.8.2；拓展：4.8.3 \\
5.1 & 后端服务与REST接口概述 & 核心：5.1.1、5.1.2 \\
5.2 & 第一个 HTTP 接口：从固定数据到契约 & 核心：5.2.1、5.2.2、5.2.3 \\
5.3 & Spring Boot 工程是怎样装起来的 & 核心：5.3.1、5.3.2、5.3.3；拓展：5.3.4 \\
5.4 & Jakarta Persistence与事务管理 & 核心：5.4.1、5.4.2、5.4.4、5.4.5、5.4.6；指导实践：5.4.3、5.4.7；拓展：5.4.8、5.4.9、5.4.10、5.4.11 \\
5.5 & 服务层与异常处理 & 核心：5.5.1、5.5.2；指导实践：5.5.3、5.5.4；拓展：5.5.5 \\
5.6 & Spring Security 6与JWT & 核心：5.6.1；指导实践：5.6.2；拓展：5.6.3 \\
5.7 & 事务事件与消息队列 & 核心：5.7.1；拓展：5.7.2、5.7.3、5.7.4、5.7.5、5.7.6。本节引言和5.7.1讲清两个异步概念：“提交之后再做”，以及同步调用为什么会把下游的故障传给上游。5.7.2—5.7.6是 Kafka 生产者、消费者幂等、重试与死信、事务发件箱和分区顺序的完整实现，属于选学，用于课程设计或自学，运行它们需要配套工程的 Compose 环境 \\
5.8 & 性能诊断与可观测性专题 & 拓展 \\
5.9 & 测试策略与质量门禁 & 指导实践：5.9.1、5.9.2；拓展：5.9.3 \\
6.1 & WebGL2与Three.js三维渲染基础 & 核心：6.1.1、6.1.2、6.1.3、6.1.4、6.1.5；拓展：6.1.6 \\
6.2 & GIS服务、坐标与二三维集成 & 核心：6.2.1、6.2.2、6.2.3；指导实践：6.2.4；拓展：6.2.5 \\
6.3 & 倾斜摄影与BIM模型构建 & 拓展：6.3.1、6.3.2、6.3.3 \\
6.4 & 数字孪生水利平台架构概览 & 拓展：6.4.1、6.4.2、6.4.3、6.4.4、6.4.5、6.4.6 \\
7.1 & 面向可视化的数据准备 & 核心：7.1.1、7.1.2；指导实践：7.1.3；拓展：7.1.4、7.1.5 \\
7.2 & 图表与实时联动 & 核心：7.2.1、7.2.3、7.2.4、7.2.5；指导实践：7.2.2；拓展：7.2.6、7.2.7、7.2.8、7.2.9 \\
7.3 & 三维场景中的监测点绘制 & 核心：7.3.2；指导实践：7.3.1；拓展：7.3.3 \\
7.4 & 监测点互动与拾取技术 & 核心：7.4.1、7.4.2；指导实践：7.4.4；拓展：7.4.3 \\
8.1 & 需求分析与业务闭环 & 核心：8.1.1、8.1.2、8.1.3 \\
8.2 & 总体架构与三维场景设计 & 核心：8.2.1；指导实践：8.2.3；拓展：8.2.2 \\
8.3 & 数据模型、前端展示与后端处理 & 核心：8.3.1、8.3.2；指导实践：8.3.3、8.3.4、8.3.5、8.3.6；拓展：8.3.7、8.3.8 \\
8.4 & 智能分析、预警与业务闭环 & 核心：8.4.1、8.4.2；拓展：8.4.3 \\
8.5 & 数字孪生水利平台案例 & 拓展：8.5.1—8.5.14 \\
8.6 & 系统部署、运维与验收 & 指导实践：8.6.1、8.6.2、8.6.3 \\
9.1 & 全书回顾 & 核心 \\
9.2 & 技术发展趋势与展望 & 核心：9.2.3；拓展：9.2.1、9.2.2 \\
9.3 & 对读者的建议与后续学习指引 & 核心：9.3.1；拓展：9.3.2、9.3.3 \\
\bottomrule
\end{longtable}
% END learning-layers
